% Shared preamble for the unofficial TMA4432 practice exams C, D and E.
% Build the exam:      pdflatex -jobname=examC examC.tex
% Build the solutions: pdflatex -jobname=examC-solutions "\def\withsol{}% Shared preamble for the unofficial TMA4432 practice exams C, D and E.
% Build the exam:      pdflatex -jobname=examC examC.tex
% Build the solutions: pdflatex -jobname=examC-solutions "\def\withsol{}\input{examC.tex}"
\documentclass[11pt,a4paper]{article}
\usepackage[T1]{fontenc}
\usepackage[utf8]{inputenc}
\usepackage{lmodern}
\usepackage{amsmath,amssymb,mathtools,bm}
\usepackage[margin=2.3cm,top=2.2cm,bottom=2.2cm,headsep=0.5cm]{geometry}
\usepackage{enumitem}
\usepackage{xcolor}
\usepackage[most]{tcolorbox}
\usepackage{fancyhdr}
\usepackage{lastpage}
\usepackage{needspace}
\usepackage{comment}
\usepackage{array}

\setlength{\parindent}{0pt}
\setlength{\parskip}{0.4em}
\AtBeginDocument{%
  \setlength{\abovedisplayskip}{6pt plus 2pt minus 2pt}%
  \setlength{\belowdisplayskip}{6pt plus 2pt minus 2pt}%
  \setlength{\abovedisplayshortskip}{3pt plus 2pt}%
  \setlength{\belowdisplayshortskip}{4pt plus 2pt minus 1pt}}

% ---------- notation (follows the formula sheet) ----------
\newcommand{\R}{\mathbb{R}}
\newcommand{\C}{\mathbb{C}}
\newcommand{\Z}{\mathbb{Z}}
\newcommand{\N}{\mathbb{N}}
\newcommand{\T}{^{\mathsf{T}}}
\newcommand{\vb}[1]{\mathbf{#1}}
\newcommand{\ten}[1]{\mathcal{#1}}
\newcommand{\dd}{\,\mathrm{d}}
\newcommand{\e}{\mathrm{e}}
\newcommand{\iu}{\mathrm{i}}
\newcommand{\F}{\mathcal{F}}
\newcommand{\norm}[1]{\lVert #1\rVert}
\newcommand{\abs}[1]{\lvert #1\rvert}
\newcommand{\ip}[2]{\langle #1,#2\rangle}
\newcommand{\op}{\circledcirc}          % outer product as in the lecture
\newcommand{\kr}{\odot}                 % Khatri--Rao
\newcommand{\had}{\ast}                 % Hadamard
\newcommand{\Hess}{\operatorname{Hess}}
\newcommand{\rank}{\operatorname{rank}}
\newcommand{\spn}{\operatorname{span}}
\DeclareMathOperator{\sinc}{sinc}

% ---------- points and boxes ----------
% points at the end of a paragraph, always flush right (TeXbook trick)
\newcommand{\pts}[1]{{\unskip\nobreak\hfil\penalty50\hskip1.5em\hbox{}\nobreak\hfil\textit{(#1\,P.)}\parfillskip=0pt\finalhyphendemerits=0\par}}
% points inside a display, placed like an equation tag
\newcommand{\ptd}[1]{\tag*{\textit{(#1\,P.)}}}
\newcommand{\ptsin}[1]{\textit{(#1\,P.)}}

\newtcolorbox{infobox}{colback=gray!8,colframe=gray!45,arc=2pt,boxrule=0.5pt,left=6pt,right=6pt,top=4pt,bottom=4pt}
\newtcolorbox{trap}{enhanced,colback=red!4,colframe=red!55!black,arc=2pt,boxrule=0.5pt,
  title={\textbf{Typical trap}},fonttitle=\small\bfseries,fontupper=\small,left=6pt,right=6pt,top=2pt,bottom=2pt,before skip=8pt,after skip=6pt}
\newtcolorbox{heart}{enhanced,colback=blue!4,colframe=blue!45!black,arc=2pt,boxrule=0.5pt,
  title={\textbf{Not on the formula sheet -- know by heart}},fonttitle=\small\bfseries,fontupper=\small,left=6pt,right=6pt,top=2pt,bottom=2pt,before skip=8pt,after skip=6pt}

% ---------- problems ----------
\newcounter{problem}
\newcommand{\problem}[2]{%
  \ifnum\value{problem}>0
    \ifdefined\withsol\par\vspace{1.4em}\Needspace{0.18\textheight}{\color{gray!60}\hrule height 0.6pt}\vspace{0.9em}%
    \else\par\bigskip\Needspace{10\baselineskip}\fi
  \fi
  \stepcounter{problem}%
  {\large\textbf{Problem \theproblem.\ (#1, #2 points)}}\par\smallskip}

\setlist[enumerate,1]{label=\textbf{\alph*)},leftmargin=2.2em,itemsep=0.35em,topsep=0.3em}
\setlist[itemize]{itemsep=0.15em,topsep=0.2em}
\newcommand{\sol}[1]{\par\medskip\textbf{#1)}~}

\ifdefined\withsol
  \specialcomment{solution}{\par\medskip\noindent{\large\textbf{Solution.}}\par\nopagebreak}{\par}
\else
  \excludecomment{solution}
\fi

% ---------- header and title ----------
\newcommand{\examletter}{X}
\newcommand{\makeheader}{%
  \pagestyle{fancy}\fancyhf{}
  \renewcommand{\headrulewidth}{0.4pt}
  \fancyhead[L]{\small TMA4432 Mathematics 4C -- Practice exam \examletter\ifdefined\withsol\ (with solutions)\fi}
  \fancyhead[R]{\small Page \thepage\ of \pageref{LastPage}}
  \thispagestyle{fancy}
  \begin{center}
    {\LARGE\textbf{Practice exam \examletter}}\\[4pt]
    {\large TMA4432 Mathematics 4C, autumn 2026\ifdefined\withsol\ -- solutions\fi}
  \end{center}
  \begin{infobox}
  \textbf{Time:} 4 hours. \textbf{Points:} 100. \textbf{Aids (category D):} the official formula sheet and a simple calculator, nothing else.\\
  All answers have to be justified, with enough detail to see how they were obtained.\\[3pt]
  \begin{tabular}{@{}l@{\hspace{1.2em}}*{5}{p{1.4cm}}l@{}}
  Grade & A & B & C & D & E & F\\
  Points & 100--89 & 88--77 & 76--65 & 64--53 & 52--41 & 40 and less
  \end{tabular}\\[3pt]
  {\small Unofficial practice exam, built from the syllabus (lectures 1--12 and the summary lecture), the exercise sheets, the test exam and Bergmann's Math 4N exams 2022--2025. Notation follows the formula sheet.}
  \end{infobox}\par\medskip}

\renewcommand{\examletter}{C}
\begin{document}
\makeheader

% =====================================================================
\problem{SVD and Moore--Penrose inverse}{10}
Consider the matrix
\[
A=\begin{pmatrix}1&0&1\\0&1&1\end{pmatrix}\in\R^{2\times 3}.
\]
\begin{enumerate}
\item Compute the singular values of $A$ and a reduced singular value decomposition $A=U\Sigma V\T$ with $\Sigma\in\R^{2\times2}$.
\item Compute the Moore--Penrose inverse $A^\dagger$.
\item Consider the linear system $A\vb x=\vb b$ with $\vb b=(1,2)\T$. Compute $\vb x^\dagger=A^\dagger\vb b$ and show that $\tilde{\vb x}=(1,2,0)\T$ also solves the system. Describe all solutions and explain why $\vb x^\dagger$ is the one of smallest norm.
\item State the best rank-1 approximation $A_1$ of $A$ and the errors $\norm{A-A_1}_2$ and $\norm{A-A_1}_{\mathrm F}$.
\end{enumerate}

\begin{solution}
\sol{a} $A$ has fewer rows than columns, so we use the smaller matrix $AA\T\in\R^{2\times2}$:
\[
AA\T=\begin{pmatrix}2&1\\1&2\end{pmatrix},\qquad \det(AA\T-\lambda I)=(2-\lambda)^2-1=(\lambda-3)(\lambda-1).
\]
Hence $\sigma_1=\sqrt3$, $\sigma_2=1$. \pts{1}

Eigenvectors of $AA\T$: $AA\T-3I=\begin{psmallmatrix}-1&1\\1&-1\end{psmallmatrix}$ gives $\vb u_1=\frac1{\sqrt2}(1,1)\T$, and $AA\T-I=\begin{psmallmatrix}1&1\\1&1\end{psmallmatrix}$ gives $\vb u_2=\frac1{\sqrt2}(1,-1)\T$. \pts{1}

The right singular vectors follow from $A\T\vb u_j=\sigma_j\vb v_j$:
\[
\vb v_1=\frac{1}{\sqrt3}\cdot\frac{1}{\sqrt2}\begin{pmatrix}1\\1\\2\end{pmatrix}=\frac1{\sqrt6}\begin{pmatrix}1\\1\\2\end{pmatrix},\qquad
\vb v_2=\frac{1}{1}\cdot\frac{1}{\sqrt2}\begin{pmatrix}1\\-1\\0\end{pmatrix}.\ptd{1}
\]

The reduced SVD is
\[
A=\underbrace{\frac1{\sqrt2}\begin{pmatrix}1&1\\1&-1\end{pmatrix}}_{U}
\underbrace{\begin{pmatrix}\sqrt3&0\\0&1\end{pmatrix}}_{\Sigma}
\underbrace{\begin{pmatrix}\frac1{\sqrt6}&\frac1{\sqrt6}&\frac2{\sqrt6}\\[2pt]\frac1{\sqrt2}&-\frac1{\sqrt2}&0\end{pmatrix}}_{V\T}.\ptd{1}
\]

\sol{b} By definition $A^\dagger=V\Sigma^{-1}U\T=\frac1{\sigma_1}\vb v_1\vb u_1\T+\frac1{\sigma_2}\vb v_2\vb u_2\T$, with the two terms
\[
\frac{1}{\sqrt3\sqrt6\sqrt2}\begin{pmatrix}1&1\\1&1\\2&2\end{pmatrix}=\frac16\begin{pmatrix}1&1\\1&1\\2&2\end{pmatrix},\qquad
\frac12\begin{pmatrix}1&-1\\-1&1\\0&0\end{pmatrix},
\qquad\text{so}\qquad
A^\dagger=\frac13\begin{pmatrix}2&-1\\-1&2\\1&1\end{pmatrix}.\ptd{1}
\]
Check: $A$ has full row rank, so $AA^\dagger=I_{2\times2}$ must hold, and indeed $AA^\dagger=\frac13\begin{psmallmatrix}2+1&-1+1\\-1+1&2+1\end{psmallmatrix}=I$. \pts{1}

\sol{c} $\vb x^\dagger=A^\dagger\vb b=\frac13(2-2,\,-1+4,\,1+2)\T=(0,1,1)\T$, and $A\vb x^\dagger=(1,2)\T$. Also $A\tilde{\vb x}=(1+0,\,2+0)\T=\vb b$. \pts{1}

The kernel is given by $x_1+x_3=0$, $x_2+x_3=0$, so $\ker A=\spn\{\vb n\}$ with $\vb n=(1,1,-1)\T$. All solutions are $\vb x=\vb x^\dagger+s\vb n$, $s\in\R$; $\tilde{\vb x}$ is the case $s=1$. Since $\ip{\vb x^\dagger}{\vb n}=0+1-1=0$,
\[
\norm{\vb x^\dagger+s\vb n}_2^2=\norm{\vb x^\dagger}_2^2+s^2\norm{\vb n}_2^2=2+3s^2\ \ge\ 2,
\]
with equality only for $s=0$. So $\norm{\vb x^\dagger}_2^2=2<5=\norm{\tilde{\vb x}}_2^2$. The Moore--Penrose solution lies in $\operatorname{range}(A\T)=\spn\{\vb v_1,\vb v_2\}$, which is orthogonal to $\ker A$. \pts{1}

\sol{d} The best rank-1 approximation is the truncated SVD
\[
A_1=\sigma_1\vb u_1\vb v_1\T=\sqrt3\cdot\frac1{\sqrt2}\cdot\frac1{\sqrt6}\begin{pmatrix}1&1&2\\1&1&2\end{pmatrix}=\frac12\begin{pmatrix}1&1&2\\1&1&2\end{pmatrix}.\ptd{1}
\]
The errors are given by the dropped singular value: $\norm{A-A_1}_2=\sigma_2=1$ and $\norm{A-A_1}_{\mathrm F}=\big(\sigma_2^2\big)^{1/2}=1$. (Check: $A-A_1=\frac12\begin{psmallmatrix}1&-1&0\\-1&1&0\end{psmallmatrix}$.) \pts{1}

\begin{trap}
\begin{itemize}
\item For a wide matrix ($m<n$) use $AA\T$, then get $\vb v_j=A\T\vb u_j/\sigma_j$. Computing the eigenvectors of $A\T A$ separately can give signs that do not match.
\item $A^\dagger$ has the transposed size, here $3\times2$. Check with $AA^\dagger=I$; $A^\dagger A$ is only a projection (rank 2).
\item ``Smallest norm'' refers to $\norm{\vb x}_2$; the residual is zero for every solution here.
\end{itemize}
\end{trap}
\begin{heart}
$A^\dagger=V\Sigma^{-1}U\T$. For full row rank $A^\dagger=A\T(AA\T)^{-1}$ and $\vb x^\dagger=A^\dagger\vb b$ is the solution of minimal norm; for full column rank it is the least-squares solution. Eckart--Young: $\norm{A-A_k}_2=\sigma_{k+1}$, $\norm{A-A_k}_{\mathrm F}=\big(\sum_{j>k}\sigma_j^2\big)^{1/2}$.
\end{heart}
\end{solution}

% =====================================================================
\problem{Matricization and mode products}{8}
Let $\ten X\in\R^{2\times2\times2}$ be given by its two frontal slices
\[
X_{::1}=\begin{pmatrix}1&2\\0&1\end{pmatrix},\qquad X_{::2}=\begin{pmatrix}2&4\\0&2\end{pmatrix}.
\]
\begin{enumerate}
\item State the three matricizations $X_{(1)}$, $X_{(2)}$ and $X_{(3)}$.
\item Compute the 1-mode product $\ten Y=\ten X\times_1 N$ with $N=\begin{pmatrix}1&-2\\0&1\end{pmatrix}$ and state its frontal slices.
\item Write $\ten X$ as a sum of two rank-1 tensors and determine the rank of $\ten X$.
\item Which matricization has rank 1? What does this mean for a Tucker decomposition of $\ten X$ with sizes $(r_1,r_2,r_3)=(2,2,1)$ computed by the HOSVD? Compute the factor matrix $C$.
\end{enumerate}

\begin{solution}
\sol{a} The column index runs through the remaining modes ``first index first'':
\[
X_{(1)}=\big(X_{::1}\ X_{::2}\big)=\begin{pmatrix}1&2&2&4\\0&1&0&2\end{pmatrix},\qquad
X_{(2)}=\big(X_{::1}\T\ X_{::2}\T\big)=\begin{pmatrix}1&0&2&0\\2&1&4&2\end{pmatrix},\ptd{2}
\]
\[
X_{(3)}=\begin{pmatrix}\operatorname{vec}(X_{::1})\T\\ \operatorname{vec}(X_{::2})\T\end{pmatrix}=\begin{pmatrix}1&0&2&1\\2&0&4&2\end{pmatrix}.
\]

\sol{b} With $Y_{(1)}=NX_{(1)}$:
\[
Y_{(1)}=\begin{pmatrix}1&-2\\0&1\end{pmatrix}\begin{pmatrix}1&2&2&4\\0&1&0&2\end{pmatrix}=\begin{pmatrix}1&0&2&0\\0&1&0&2\end{pmatrix},
\qquad\text{so}\qquad Y_{::1}=I_{2\times2},\quad Y_{::2}=2I_{2\times2}.
\]
Each frontal slice is multiplied from the left, $Y_{::k}=NX_{::k}$; here $N=X_{::1}^{-1}$. \pts{2}

\sol{c} Since $X_{::2}=2X_{::1}$ we have $x_{ijk}=m_{ij}c_k$ with $M=X_{::1}$ and $\vb c=(1,2)\T$. Splitting $M$ into its rows, $M=\vb e_1(1,2)+\vb e_2(0,1)$, gives
\[
\ten X=\vb e_1\op\begin{pmatrix}1\\2\end{pmatrix}\op\vb c\;+\;\vb e_2\op\begin{pmatrix}0\\1\end{pmatrix}\op\vb c,
\]
so $\rank\ten X\le2$. \pts{1}
For a sum of $r$ rank-1 tensors, $X_{(1)}=A(C\kr B)\T$ has rank at most $r$. The first two columns of $X_{(1)}$, $(1,0)\T$ and $(2,1)\T$, are linearly independent, so $\rank X_{(1)}=2$ and therefore $\rank\ten X=2$. \pts{1}

\sol{d} In $X_{(3)}$ the second row is twice the first, so $\rank X_{(3)}=1$: every mode-3 fiber $\vb x_{ij:}=m_{ij}\vb c$ is a multiple of $\vb c$. \pts{1}

The HOSVD takes $C$ as the leading left singular vector of $X_{(3)}=\vb c\operatorname{vec}(M)\T$, i.e.
$C=\vb c/\norm{\vb c}_2=\frac1{\sqrt5}(1,2)\T$ (singular value $\norm{\vb c}_2\norm{M}_{\mathrm F}=\sqrt5\cdot\sqrt6=\sqrt{30}$).
Projecting the mode-3 fibers onto $\spn\{\vb c\}$ changes nothing, and for $r_1=r_2=2$ the factors $A,B$ are orthogonal $2\times2$ matrices. Hence
$\ten G\times_1A\times_2B\times_3C=\ten X\times_1AA\T\times_2BB\T\times_3CC\T=\ten X$: the Tucker decomposition is exact with only $2\cdot2\cdot1=4$ core entries instead of 8. (With $A=B=I$ the core is $\ten G=\ten X\times_3C\T=\sqrt5\,M$.) \pts{1}

\begin{trap}
\begin{itemize}
\item Column order of $X_{(k)}$: the first remaining index runs fastest. $X_{(2)}$ consists of the \emph{transposed} slices.
\item $\ten X\times_1N$ means $Y_{(1)}=NX_{(1)}$, not $X_{(1)}N$.
\item The rank of a tensor is at least the rank of each matricization, but in general not equal to it.
\end{itemize}
\end{trap}
\begin{heart}
$\ten Y=\ten X\times_kA\iff Y_{(k)}=AX_{(k)}$. For $\ten X=[\![A,B,C]\!]$: $X_{(1)}=A(C\kr B)\T$, $X_{(2)}=B(C\kr A)\T$, $X_{(3)}=C(B\kr A)\T$.
HOSVD: $A,B,C$ are the leading left singular vectors of $X_{(1)},X_{(2)},X_{(3)}$, and $\ten G=\ten X\times_1A\T\times_2B\T\times_3C\T$.
\end{heart}
\end{solution}

% =====================================================================
\problem{A curve in space}{8}
Consider the curve $C$ given by the parametrization
\[
\vb c(t)=\begin{pmatrix}2\cos t\\2\sin t\\\sqrt5\,t\end{pmatrix},\qquad t\in\R.
\]
\begin{enumerate}
\item Compute the velocity $\vb v(t)$, the speed and the acceleration $\vb a(t)$. Is the parametrization smooth?
\item Determine the tangent line and the unit tangent vector at the point $\vb p=\vb c(\pi/2)$.
\item Compute the arc length from $\vb c(0)$ to $\vb c(t)$, $t>0$, and state the arc length parametrization of $C$.
\item Consider $\vb d(t)=(t^3,t^6)\T$, $t\in\R$. Is the parametrization $\vb d$ smooth? Is the curve $D=\{\vb d(t):t\in\R\}$ smooth?
\end{enumerate}

\begin{solution}
\sol{a}
\[
\vb v(t)=\vb c'(t)=\begin{pmatrix}-2\sin t\\2\cos t\\\sqrt5\end{pmatrix},\qquad
\norm{\vb v(t)}_2=\sqrt{4\sin^2t+4\cos^2t+5}=3,\qquad
\vb a(t)=\begin{pmatrix}-2\cos t\\-2\sin t\\0\end{pmatrix}.\ptd{1}
\]
$\vb c$ is differentiable, $\vb c'(t)\neq\vb 0$ for all $t$, and $\vb c$ is injective since the third component $\sqrt5\,t$ is strictly increasing. So the parametrization is smooth. \pts{1}

\sol{b} $\vb p=\vb c(\pi/2)=(0,2,\sqrt5\pi/2)\T$ and $\vb c'(\pi/2)=(-2,0,\sqrt5)\T$. The tangent line is
\[
L(s)=\begin{pmatrix}0\\2\\\sqrt5\pi/2\end{pmatrix}+s\begin{pmatrix}-2\\0\\\sqrt5\end{pmatrix},\quad s\in\R,\ptd{1}
\]
and the unit tangent vector is $\frac13(-2,0,\sqrt5)\T$. \pts{1}

\sol{c} $\ell(t)=\int_0^t\norm{\vb c'(\tau)}_2\dd\tau=\int_0^t3\dd\tau=3t$. \pts{1}
So $s=3t$, i.e.\ $t=s/3$, and the arc length parametrization is
\[
\tilde{\vb c}(s)=\vb c(s/3)=\begin{pmatrix}2\cos(s/3)\\2\sin(s/3)\\\sqrt5\,s/3\end{pmatrix},\qquad \norm{\tilde{\vb c}'(s)}_2=\tfrac13\cdot3=1.\ptd{1}
\]

\sol{d} $\vb d'(t)=(3t^2,6t^5)\T$ vanishes at $t=0$, so the parametrization $\vb d$ is not smooth in $t=0$. \pts{1}
But $x=t^3$ runs through all of $\R$ and $y=t^6=x^2$, so $D=\{(x,x^2)\T:x\in\R\}$ is a parabola. It has the parametrization $\vb e(s)=(s,s^2)\T$ with $\vb e'(s)=(1,2s)\T\neq\vb 0$, and $\vb e$ is injective. Hence the curve $D$ is smooth: a curve is smooth if it has \emph{some} smooth parametrization. \pts{1}

\begin{trap}
\begin{itemize}
\item Distinguish a smooth \emph{parametrization} from a smooth \emph{curve}: $\vb d$ is not smooth, the curve $D$ is.
\item The speed is the norm of the velocity, a number; the tangent line needs a point \emph{and} a direction.
\item An arc length parametrization has unit speed. Check it.
\end{itemize}
\end{trap}
\begin{heart}
Arc length $\ell=\int_a^b\norm{\vb c'(t)}_2\dd t$; tangent line $\vb c(t_0)+s\,\vb c'(t_0)$; unit tangent $\vb c'(t_0)/\norm{\vb c'(t_0)}_2$; smooth $=$ differentiable, injective and $\vb c'(t)\neq\vb 0$.
\end{heart}
\end{solution}

% =====================================================================
\problem{Newton's method for a system}{7}
Consider $\vb F\colon\R^2\to\R^2$,
\[
\vb F(x,y)=\begin{pmatrix}x^3-y\\y^3-x\end{pmatrix}.
\]
\begin{enumerate}
\item Compute the Jacobi matrix $\vb F'(x,y)$ and all real zeros of $\vb F$.
\item Perform one step of Newton's method starting from $(x_0,y_0)=(2,1)$.
\item For which starting points $(a,a)$ on the diagonal can the first Newton step not be computed? Which convergence speed do you expect close to the zero $(1,1)$, and why?
\end{enumerate}

\begin{solution}
\sol{a}
\[
\vb F'(x,y)=\begin{pmatrix}3x^2&-1\\-1&3y^2\end{pmatrix},\qquad \det\vb F'(x,y)=9x^2y^2-1.\ptd{1}
\]
Zeros: $y=x^3$ and $x=y^3=x^9$, so $x(x^8-1)=0$ and $x\in\{0,1,-1\}$. The zeros are $(0,0)$, $(1,1)$ and $(-1,-1)$. \pts{1}

\sol{b} $\vb F(2,1)=(7,-1)\T$ and $\vb F'(2,1)=\begin{psmallmatrix}12&-1\\-1&3\end{psmallmatrix}$ with determinant $35$. \pts{1}
We solve $\vb F'(2,1)\Delta=-\vb F(2,1)$:
\[
\Delta=\frac1{35}\begin{pmatrix}3&1\\1&12\end{pmatrix}\begin{pmatrix}-7\\1\end{pmatrix}=\frac1{35}\begin{pmatrix}-20\\5\end{pmatrix}=\begin{pmatrix}-4/7\\1/7\end{pmatrix}.\ptd{1}
\]
Hence $(x_1,y_1)=(2,1)+\Delta=\big(\tfrac{10}7,\tfrac87\big)\approx(1.429,\,1.143)$. \pts{1}

\sol{c} On the diagonal $\det\vb F'(a,a)=9a^4-1=0$ for $a=\pm1/\sqrt3$. There the linear system for $\Delta$ is singular and the step is not defined. \pts{1}
At the zero, $\det\vb F'(1,1)=8\neq0$, so the Jacobi matrix is invertible there. Then Newton's method converges quadratically for starting points close enough to $(1,1)$. From $(2,1)$ the errors are about $0.45,\ 0.13,\ 0.016,\ 2.9\cdot10^{-4},\ 9\cdot10^{-8}$: the number of correct digits roughly doubles. \pts{1}

\begin{trap}
\begin{itemize}
\item The step solves $\vb F'(\vb x^{(k)})\Delta=-\vb F(\vb x^{(k)})$. The minus sign is the most common error.
\item Convergence is only local. Newton does not necessarily find the nearest zero: from $(\frac12,\frac12)$ one step lands exactly on $(-1,-1)$, not on $(0,0)$.
\end{itemize}
\end{trap}
\begin{heart}
$\begin{psmallmatrix}a&b\\c&d\end{psmallmatrix}^{-1}=\frac1{ad-bc}\begin{psmallmatrix}d&-b\\-c&a\end{psmallmatrix}$. Newton converges quadratically if $\vb F'(\vb x^*)$ is invertible and $\vb x^{(0)}$ is close enough to $\vb x^*$. (The Newton step itself is on the formula sheet.)
\end{heart}
\end{solution}

% =====================================================================
\problem{Optimization}{10}
Consider $f\colon\R^2\to\R$, $f(x,y)=x^3-3xy+y^3$.
\begin{enumerate}
\item Compute the gradient $\nabla f$ and all stationary points of $f$.
\item Compute the Hessian $\Hess f$. Show that $(0,0)$ is a saddle point by considering $f$ along the lines $y=x$ and $y=-x$. Does $f$ have a global minimizer?
\item Perform one gradient descent step from $\vb x^{(0)}=(1,0)\T$ with step size $s=\frac13$ and compare $f(\vb x^{(1)})$ with $f(\vb x^{(0)})$. Then halve the step size until the function value decreases (backtracking). Which step size and which $\vb x^{(1)}$ do you get?
\item Perform one step of Newton's method for the minimization of $f$ (Newton applied to $\nabla f(\vb x)=\vb 0$), starting from $(2,2)\T$. Why can Newton's method on its own not guarantee that the limit is a minimizer?
\end{enumerate}

\begin{solution}
\sol{a} $\nabla f(x,y)=(3x^2-3y,\;3y^2-3x)\T$. \pts{1}
$\nabla f=\vb 0\iff y=x^2$ and $x=y^2=x^4$, so $x(x^3-1)=0$. The stationary points are $(0,0)$ and $(1,1)$. \pts{1}

\sol{b}
\[
\Hess f(x,y)=\begin{pmatrix}6x&-3\\-3&6y\end{pmatrix}.\ptd{1}
\]
Along $y=x$: $f(t,t)=2t^3-3t^2=t^2(2t-3)<0$ for small $t\neq0$. Along $y=-x$: $f(t,-t)=t^3+3t^2-t^3=3t^2>0$ for $t\neq0$. Since $f(0,0)=0$, $f$ takes both smaller and larger values arbitrarily close to $(0,0)$, so $(0,0)$ is a saddle point. \pts{1}
$f(1,1)=-1$, but $f(x,0)=x^3\to-\infty$ for $x\to-\infty$ (e.g.\ $f(-2,0)=-8<-1$), so $f$ has no global minimizer; $(1,1)$ can at most be a local one. \pts{1}

\sol{c} $\nabla f(1,0)=(3,-3)\T$ and $f(1,0)=1$. With $s=\frac13$:
$\vb x^{(1)}=(1,0)\T-\frac13(3,-3)\T=(0,1)\T$ and $f(0,1)=1$: no decrease. \pts{1}
Halving gives $s=\frac16$: $\vb x^{(1)}=(1,0)\T-\frac16(3,-3)\T=(\frac12,\frac12)\T$ with $f(\frac12,\frac12)=\frac18-\frac34+\frac18=-\frac12<1$. Accept $s=\frac16$. \pts{1}

\sol{d} $\nabla f(2,2)=(6,6)\T$ and $\Hess f(2,2)=\begin{psmallmatrix}12&-3\\-3&12\end{psmallmatrix}$. \pts{1}
Solve $\Hess f(2,2)\Delta=-\nabla f(2,2)$. By symmetry $\Delta=(d,d)\T$ with $(12-3)d=-6$, so $d=-\frac23$ and $\vb x^{(1)}=(\frac43,\frac43)\T$. (The next step gives $\frac{16}{15}(1,1)\T$, converging to $(1,1)$.) \pts{1}
Newton's method solves $\nabla f(\vb x)=\vb 0$ and converges to stationary points of any type; started close to $(0,0)$ it converges to the saddle point. Whether a limit is a minimizer has to be checked separately. Here $(1,1)$ is a local minimizer, since $f(1+h,1+k)=-1+3\big((h-\frac k2)^2+\frac34k^2\big)+h^3+k^3>-1$ for small $(h,k)\neq(0,0)$. \pts{1}

\begin{trap}
\begin{itemize}
\item $\nabla f(\vb x^*)=\vb 0$ is necessary for a local minimum, not sufficient: $(0,0)$ is a saddle point. To show a saddle point, two directions with opposite behaviour are enough.
\item ``Local minimizer'' and ``global minimizer'' are different claims; here $f$ is unbounded below.
\item A fixed step size can fail to decrease $f$; that is the reason for backtracking.
\end{itemize}
\end{trap}
\begin{heart}
Gradient descent $\vb x^{(k+1)}=\vb x^{(k)}-s_k\nabla f(\vb x^{(k)})$; backtracking: try $s=\beta^\ell u$, $\ell=0,1,\dots$, until the decrease is large enough (Armijo). Newton for optimization: $\Hess f(\vb x^{(k)})\Delta^{(k)}=-\nabla f(\vb x^{(k)})$. A local extremum in the interior has $\nabla f=\vb 0$, but not the reverse.
\end{heart}
\end{solution}

% =====================================================================
\problem{Fourier series}{12}
Let $f(x)=x^2$ for $-\pi\le x<\pi$, extended $2\pi$-periodically.
\begin{enumerate}
\item Compute the complex Fourier coefficients $c_k$ of $f$, and from them the real coefficients $a_0,a_n,b_n$.
\item To which value does the Fourier series of $f$ converge at $x=\pi$? Use this to compute $\sum_{n=1}^\infty\frac1{n^2}$.
\item Use Parseval's identity to compute $\sum_{k=1}^\infty\frac1{k^4}$.
\item Use the derivative rule for Fourier coefficients to determine the complex Fourier coefficients of $g(x)=x$, $-\pi\le x<\pi$ ($2\pi$-periodic), without computing new integrals. Why may the rule be applied here?
\item Compute the convolution $h=f\ast s$ with $s(x)=\cos(2x)$.
\end{enumerate}

\begin{solution}
\sol{a} $c_0=\frac1{2\pi}\int_{-\pi}^\pi x^2\dd x=\frac{\pi^2}3$. \pts{1}
For $k\neq0$ we integrate by parts twice, using $\e^{\pm\iu k\pi}=(-1)^k$:
\begin{align*}
\int_{-\pi}^{\pi}x^2\e^{-\iu kx}\dd x&=\underbrace{\Big[\frac{x^2\e^{-\iu kx}}{-\iu k}\Big]_{-\pi}^{\pi}}_{=0}+\frac2{\iu k}\int_{-\pi}^\pi x\e^{-\iu kx}\dd x,\\
\int_{-\pi}^\pi x\e^{-\iu kx}\dd x&=\Big[\frac{x\e^{-\iu kx}}{-\iu k}\Big]_{-\pi}^{\pi}+\underbrace{\frac1{\iu k}\int_{-\pi}^\pi\e^{-\iu kx}\dd x}_{=0}=\frac{2\pi(-1)^k}{-\iu k}=\frac{2\pi\iu(-1)^k}{k}.
\end{align*}
Hence $\int_{-\pi}^\pi x^2\e^{-\iu kx}\dd x=\frac{4\pi(-1)^k}{k^2}$ and $c_k=\frac{2(-1)^k}{k^2}$, $k\neq0$. \pts{1}
Real coefficients: $a_0=c_0=\frac{\pi^2}3$, $a_n=c_n+c_{-n}=\frac{4(-1)^n}{n^2}$, $b_n=\iu(c_n-c_{-n})=0$ ($f$ is even). \pts{1}

\sol{b} The periodic extension is continuous ($f(-\pi)=f(\pi)=\pi^2$, no jump) and piecewise continuously differentiable, so the series converges to $f(\pi)=\pi^2$ at $x=\pi$. \pts{1}
Since $(-1)^n\cos(n\pi)=1$,
\[
\pi^2=\frac{\pi^2}{3}+\sum_{n=1}^\infty\frac{4(-1)^n}{n^2}\cos(n\pi)=\frac{\pi^2}{3}+4\sum_{n=1}^\infty\frac1{n^2}
\quad\Longrightarrow\quad\sum_{n=1}^\infty\frac1{n^2}=\frac{\pi^2}6.\ptd{1}
\]

\sol{c} Parseval: $\frac1{2\pi}\int_{-\pi}^\pi\abs{f(x)}^2\dd x=\sum_{k\in\Z}\abs{c_k}^2$. The left-hand side is $\frac1{2\pi}\cdot\frac{2\pi^5}{5}=\frac{\pi^4}5$. \pts{1}
The right-hand side is $\frac{\pi^4}9+2\sum_{k=1}^\infty\frac4{k^4}$ (the terms $k$ and $-k$ are equal). \pts{1}
Hence
\[
\sum_{k=1}^\infty\frac1{k^4}=\frac18\Big(\frac{\pi^4}5-\frac{\pi^4}9\Big)=\frac18\cdot\frac{4\pi^4}{45}=\frac{\pi^4}{90}.\ptd{1}
\]

\sol{d} The periodic extension of $f$ is continuous and $f'(x)=2x$ is piecewise continuous, so $c_k(f')=\iu k\,c_k(f)$ holds. \pts{1}
With $f'=2g$:
\[
c_k(g)=\frac12\,\iu k\,\frac{2(-1)^k}{k^2}=\frac{\iu(-1)^k}{k},\quad k\neq0,\qquad c_0(g)=\frac12\cdot0\cdot c_0(f)=0.
\]
Check: $b_n=\iu(c_n-c_{-n})=\frac{2(-1)^{n+1}}{n}$, the known sine series of $x$. \pts{1}

\sol{e} By the convolution theorem $c_k(f\ast s)=2\pi\,c_k(f)\,c_k(s)$. Since $s(x)=\frac12(\e^{2\iu x}+\e^{-2\iu x})$, only $c_{\pm2}(s)=\frac12$ are nonzero. \pts{1}
So $c_{\pm2}(h)=2\pi\cdot\frac{2}{4}\cdot\frac12=\frac\pi2$ and all other coefficients vanish:
\[
h(x)=\frac\pi2\big(\e^{2\iu x}+\e^{-2\iu x}\big)=\pi\cos(2x).\ptd{1}
\]

\begin{trap}
\begin{itemize}
\item $c_0$ has to be computed separately; the formula $2(-1)^k/k^2$ is only valid for $k\neq0$.
\item In Parseval both $k$ and $-k$ appear on the right, and the left-hand side carries the factor $\frac1{2\pi}$.
\item The derivative rule needs a \emph{continuous} periodic extension. It may not be applied to $g(x)=x$ itself (jump at $\pm\pi$).
\item In this course $(f\ast g)(x)=\int_{-\pi}^\pi f(y)g(x-y)\dd y$, without $\frac1{2\pi}$. Hence the factor $2\pi$ in $c_k(f\ast g)$.
\end{itemize}
\end{trap}
\begin{heart}
$a_n=c_n+c_{-n}$, $b_n=\iu(c_n-c_{-n})$, $c_{\pm n}=\frac12(a_n\mp\iu b_n)$. Parseval: $\frac1{2\pi}\int_{-\pi}^{\pi}\abs f^2\dd x=\sum_k\abs{c_k}^2=a_0^2+\frac12\sum_{n\ge1}(a_n^2+b_n^2)$.
$c_k(f\ast g)=2\pi c_k(f)c_k(g)$, $c_k(f')=\iu k\,c_k(f)$. The series converges to $\frac12\big(f(x^-)+f(x^+)\big)$.
\end{heart}
\end{solution}

% =====================================================================
\problem{Discrete Fourier transform}{10}
Consider the signal $\vb f=(1,2,0,-1)\T\in\R^4$.
\begin{enumerate}
\item State the Fourier matrix $\F_4$ and compute the DFT $\hat{\vb f}$.
\item Verify that $\hat f_3=\overline{\hat f_1}$ and explain why this has to hold.
\item The signal consists of the samples $f_j=p(x_j)$, $x_j=\frac{2\pi j}{4}$, of a real trigonometric polynomial $p(x)=a_0+a_1\cos x+b_1\sin x$. Determine $a_0,a_1,b_1$ from $\hat{\vb f}$ and check your result at the four sample points.
\item For $N=2^{10}=1024$, compare the number of operations needed to compute the DFT as a matrix--vector product and with the FFT.
\end{enumerate}

\begin{solution}
\sol{a} With $w_4=\e^{-2\pi\iu/4}=-\iu$,
\[
\F_4=\big(w_4^{jk}\big)_{j,k=0}^3=\begin{pmatrix}1&1&1&1\\1&-\iu&-1&\iu\\1&-1&1&-1\\1&\iu&-1&-\iu\end{pmatrix}.\ptd{1}
\]
$\hat{\vb f}=\frac14\F_4\vb f$:
\begin{align*}
\hat f_0&=\tfrac14(1+2+0-1)=\tfrac12, & \hat f_1&=\tfrac14(1-2\iu+0-\iu)=\tfrac{1-3\iu}4,\\
\hat f_2&=\tfrac14(1-2+0+1)=0, & \hat f_3&=\tfrac14(1+2\iu+0+\iu)=\tfrac{1+3\iu}4.\ptd{2}
\end{align*}

\sol{b} $\hat f_3=\frac{1+3\iu}4=\overline{\frac{1-3\iu}4}=\overline{\hat f_1}$. \pts{1}
For a real signal, $\hat f_{N-k}=\frac1N\sum_jf_j\e^{-2\pi\iu j(N-k)/N}=\frac1N\sum_jf_j\e^{2\pi\iu jk/N}=\overline{\hat f_k}$, because $\e^{-2\pi\iu j}=1$ and $f_j=\overline{f_j}$. \pts{1}

\sol{c} By the aliasing formula $\hat f_k=\sum_{\ell}c_{k+4\ell}(p)$, and $p$ only has $c_0,c_{\pm1}$. So $\hat f_0=c_0$, $\hat f_1=c_1$, $\hat f_3=\hat f_{-1}=c_{-1}$ (and $\hat f_2=0$ fits). \pts{1}
Hence $a_0=c_0=\frac12$, $a_1=c_1+c_{-1}=\frac{1-3\iu}4+\frac{1+3\iu}4=\frac12$ and $b_1=\iu(c_1-c_{-1})=\iu\cdot\frac{-6\iu}4=\frac32$. \pts{1}
So $p(x)=\frac12+\frac12\cos x+\frac32\sin x$, and $p(0)=1$, $p(\frac\pi2)=\frac12+\frac32=2$, $p(\pi)=0$, $p(\frac{3\pi}2)=\frac12-\frac32=-1$. \pts{1}

\sol{d} As a matrix--vector product the DFT needs $N^2=1\,048\,576$ multiplications, $\mathcal O(N^2)$. The FFT needs $\mathcal O(N\log_2N)$: about $\frac N2\log_2N=5120$ multiplications and $N\log_2N=10\,240$ additions. That is about 100 times fewer operations. \pts{2}

\begin{trap}
\begin{itemize}
\item The factor $\frac1N$ belongs to the DFT (formula sheet), not to the inverse DFT.
\item $\hat f_3$ belongs to the frequency $-1$, not $3$ (FFT-shift).
\item $w_N=\e^{-2\pi\iu/N}$ carries a minus sign.
\end{itemize}
\end{trap}
\begin{heart}
Aliasing formula $\hat f_k=\sum_{\ell\in\Z}c_{k+\ell N}(f)$. For real signals $\hat f_{N-k}=\overline{\hat f_k}$. DFT: $\mathcal O(N^2)$, FFT: $\mathcal O(N\log N)$.
\end{heart}
\end{solution}

% =====================================================================
\problem{Heat equation on the real line}{14}
Consider the heat equation on the real line
\[
\frac{\partial u}{\partial t}=2\frac{\partial^2u}{\partial x^2},\quad x\in\R,\ t>0,\qquad u(x,0)=\e^{-x^2},\qquad u(x,t)\to0\ \text{for } \abs x\to\infty.
\]
\begin{enumerate}
\item Classify the PDE (order, linear/nonlinear, homogeneous/inhomogeneous). Why does the superposition principle hold?
\item Apply the Fourier transform in $x$, derive the ODE for $\hat u(\omega,t)$ and solve it using the initial condition.
\item Compute $\hat u(\omega,t)$ explicitly and transform back to obtain $u(x,t)$.
\item Verify that your $u$ solves the PDE and the initial condition.
\item Compute $\max_{x\in\R}u(x,t)$ and $\int_{-\infty}^\infty u(x,t)\dd x$. Interpret the result.
\end{enumerate}

\begin{solution}
\sol{a} The PDE is of second order ($u_{xx}$), linear ($u$ and its derivatives appear only linearly) and homogeneous (no term without $u$). \pts{1}
If $u_1,u_2$ solve it, so does $\alpha u_1+\beta u_2$, because $\partial_t$ and $\partial_x^2$ are linear. \pts{1}

\sol{b} Let $\hat u(\omega,t)=\F\big(u(\cdot,t)\big)(\omega)$. Applying the derivative rule twice gives $\F(u_{xx})=(\iu\omega)^2\hat u=-\omega^2\hat u$, and $\F(u_t)=\partial_t\hat u$. \pts{1}
So for every fixed $\omega$:
\[
\frac{\partial\hat u}{\partial t}(\omega,t)=-2\omega^2\,\hat u(\omega,t).\ptd{1}
\]
Its solution is $\hat u(\omega,t)=C(\omega)\e^{-2\omega^2t}$ with $C(\omega)=\hat u(\omega,0)=\F(\e^{-x^2})(\omega)$. \pts{1}

\sol{c} The formula sheet ($a=1$) gives $\F(\e^{-x^2})(\omega)=\frac1{\sqrt2}\e^{-\omega^2/4}$. \pts{1}
So $\hat u(\omega,t)=\frac1{\sqrt2}\e^{-\omega^2/4-2\omega^2t}=\frac1{\sqrt2}\e^{-\omega^2(1+8t)/4}$. \pts{1}
Let $a=\frac1{1+8t}$. The table gives $\F(\e^{-ax^2})(\omega)=\frac1{\sqrt{2a}}\e^{-\omega^2/(4a)}=\frac1{\sqrt{2a}}\e^{-\omega^2(1+8t)/4}$, hence
$\hat u(\omega,t)=\frac1{\sqrt2}\sqrt{2a}\,\F(\e^{-ax^2})(\omega)=\sqrt a\,\F(\e^{-ax^2})(\omega)$. \pts{1}
Transforming back,
\[
u(x,t)=\sqrt a\,\e^{-ax^2}=\frac1{\sqrt{1+8t}}\exp\Big(-\frac{x^2}{1+8t}\Big).\ptd{1}
\]

\sol{d} Write $s=1+8t$, so $u=s^{-1/2}\e^{-x^2/s}$ and $\partial_t=8\,\partial_s$:
\[
u_t=8\Big(-\tfrac12s^{-3/2}+s^{-1/2}\tfrac{x^2}{s^2}\Big)\e^{-x^2/s}=s^{-5/2}\,(8x^2-4s)\,\e^{-x^2/s}.\ptd{1}
\]
\[
u_x=-2x\,s^{-3/2}\e^{-x^2/s},\qquad u_{xx}=\big(-2s^{-3/2}+4x^2s^{-5/2}\big)\e^{-x^2/s}=s^{-5/2}(4x^2-2s)\,\e^{-x^2/s},
\]
so $2u_{xx}=u_t$. \pts{1}
For $t=0$ we have $s=1$ and $u(x,0)=\e^{-x^2}$, and $u\to0$ for $\abs x\to\infty$. \pts{1}

\sol{e} $\max_xu(x,t)=u(0,t)=\frac1{\sqrt{1+8t}}\to0$ for $t\to\infty$: the peak decreases. \pts{1}
With $\int_{-\infty}^\infty\e^{-x^2/s}\dd x=\sqrt{\pi s}$ we get $\int u\dd x=s^{-1/2}\sqrt{\pi s}=\sqrt\pi$ for all $t$. The total amount of heat is conserved; it only spreads out, with a width proportional to $\sqrt{1+8t}$. \pts{1}

\begin{trap}
\begin{itemize}
\item Use the normalization of the formula sheet: $\F(\e^{-x^2})=\frac1{\sqrt2}\e^{-\omega^2/4}$ (some books get $\sqrt\pi\,\e^{-\omega^2/4}$).
\item $\F(u_{xx})=-\omega^2\hat u$; a wrong sign gives a solution that grows.
\item The ``constant'' of the ODE depends on $\omega$: $C(\omega)=\hat u(\omega,0)$.
\item Transform back by recognising a Gaussian in the table, not by integrating.
\end{itemize}
\end{trap}
\begin{heart}
$\F(f')=\iu\omega\F(f)$, $\F(f'')=-\omega^2\F(f)$, $\F(f\ast g)=\sqrt{2\pi}\,\F(f)\F(g)$.
Heat equation on $\R$: $\hat u(\omega,t)=\hat f(\omega)\e^{-c^2\omega^2t}$, then transform back.
\end{heart}
\end{solution}

% =====================================================================
\problem{Finite differences for the heat equation}{10}
Consider
\[
u_t=u_{xx},\quad 0<x<1,\ t>0,\qquad u(0,t)=u(1,t)=0,\qquad u(x,0)=f(x)=2\big(1-\abs{2x-1}\big).
\]
Use the grid $x_i=ih$, $h=\frac14$ ($M=4$), $t_n=nk$, approximations $U_i^n\approx u(x_i,t_n)$ and $\alpha=k/h^2$.
\begin{enumerate}
\item Derive the explicit scheme (forward Euler) for the inner nodes and state its approximation order in $h$ and $k$.
\item Compute $\vb U^1$ and $\vb U^2$ at the inner nodes for $k=\frac1{32}$.
\item Repeat b) with $k=\frac1{16}$. What do you observe? Explain it with the stability condition.
\item State the implicit Euler method in matrix form and compute one step with $k=\frac1{16}$. Why can this time step be used here?
\end{enumerate}

\begin{solution}
\sol{a} A forward difference in time and the central second difference in space give
\begin{align*}
u_t(x_i,t_n)&=\frac{u(x_i,t_{n+1})-u(x_i,t_n)}{k}+\mathcal O(k),\\
u_{xx}(x_i,t_n)&=\frac{u(x_{i-1},t_n)-2u(x_i,t_n)+u(x_{i+1},t_n)}{h^2}+\mathcal O(h^2).\ptd{1}
\end{align*}
Replacing $u$ by $U$:
$U_i^{n+1}=U_i^n+\alpha\big(U_{i-1}^n-2U_i^n+U_{i+1}^n\big)$, $i=1,\dots,M-1$, with $U_0^n=U_M^n=0$ and $U_i^0=f(x_i)$. The error is $\mathcal O(k+h^2)$. \pts{1}

\sol{b} $\vb U^0=\big(f(\frac14),f(\frac12),f(\frac34)\big)=(1,2,1)$ and $\alpha=\frac{1/32}{1/16}=\frac12$. \pts{1}
\[
U^1_1=1+\tfrac12(0-2+2)=1,\quad U^1_2=2+\tfrac12(1-4+1)=1,\quad U^1_3=1,\qquad \vb U^1=(1,1,1).\ptd{1}
\]
\[
U^2_1=1+\tfrac12(0-2+1)=\tfrac12,\quad U^2_2=1+\tfrac12(1-2+1)=1,\quad U^2_3=\tfrac12,\qquad \vb U^2=(\tfrac12,1,\tfrac12).\ptd{1}
\]

\sol{c} Now $\alpha=\frac1{16}\cdot16=1$:
$\vb U^1=(1+(0-2+2),\ 2+(1-4+1),\ 1)=(1,0,1)$ and $\vb U^2=(1+(0-2+0),\ 0+(1-0+1),\ -1)=(-1,2,-1)$; the next step gives $(3,-4,3)$. \pts{1}
Negative temperatures and growing oscillations: the scheme is unstable. The explicit scheme requires $\alpha=c^2k/h^2\le\frac12$, which is violated for $\alpha=1$. \pts{1}

\sol{d} A backward difference in time gives $U_i^{n+1}-\alpha\big(U_{i-1}^{n+1}-2U_i^{n+1}+U_{i+1}^{n+1}\big)=U_i^n$. In matrix form
\[
(I-\alpha B)\vb U^{n+1}=\vb U^n+\alpha\vb b^{n+1},\qquad B=\begin{pmatrix}-2&1&0\\1&-2&1\\0&1&-2\end{pmatrix},
\]
where $\vb b^{n+1}$ contains the boundary values and is $\vb 0$ here. \pts{1}
With $\alpha=1$:
\[
\begin{pmatrix}3&-1&0\\-1&3&-1\\0&-1&3\end{pmatrix}\vb U^1=\begin{pmatrix}1\\2\\1\end{pmatrix}.
\]
By symmetry $U_1=U_3=p$, $U_2=q$, with $3p-q=1$ and $-2p+3q=2$. This gives $p=\frac57$, $q=\frac87$, i.e.\ $\vb U^1=(\frac57,\frac87,\frac57)$. \pts{1}
Implicit Euler is unconditionally stable, so there is no restriction on $\alpha$. For accuracy, $k$ still has to be small (error $\mathcal O(k+h^2)$). \pts{1}

\begin{trap}
\begin{itemize}
\item $\alpha=c^2k/h^2$ with $h$ \emph{squared}.
\item Do not forget the boundary values $U_0=U_M=0$ in the first and the last inner node.
\item Implicit Euler gives a linear system, not a formula for each $U_i^{n+1}$ separately.
\end{itemize}
\end{trap}
\begin{heart}
Forward/backward difference $\mathcal O(k)$, central differences $\mathcal O(h^2)$. Explicit Euler for the heat equation is stable iff $\alpha\le\frac12$. Implicit Euler and Crank--Nicolson are unconditionally stable; Crank--Nicolson has error $\mathcal O(h^2+k^2)$.
\end{heart}
\end{solution}

% =====================================================================
\problem{Cubic splines}{11}
Let $\Delta=[0,1,2,3]$ and $(y)_+=\max\{y,0\}$.
\begin{enumerate}
\item Which of the following functions are in $S_{3,\Delta}$? Give a reason for each.
\[
f_1(x)=\abs{x-2}^3,\qquad f_2(x)=\big(x-\tfrac12\big)_+^3,\qquad f_3(x)=x^3+(x-1)_+^2.
\]
\item What is the dimension of $S_{3,\Delta}$? Why does interpolation of four values need two additional conditions?
\item Compute the natural cubic spline $s\in S_{3,\Delta}$ with $s(0)=0$, $s(1)=1$, $s(2)=0$, $s(3)=1$, i.e.\ determine $M_i=s''(x_i)$, $i=0,\dots,3$.
\item State $s$ on $[0,1]$ and compute $s(\frac12)$.
\end{enumerate}

\begin{solution}
We have to check: $s|_{[x_j,x_{j+1}]}\in\mathcal P_3$ for every interval, and $s\in C^2([0,3])$.

\sol{a} $f_1(x)=(2-x)^3$ on $[0,2]$ and $(x-2)^3$ on $[2,3]$, a cubic polynomial on each of $[0,1],[1,2],[2,3]$. At $x=2$ the values ($0$), the first derivatives ($-3(2-x)^2\to0$, $3(x-2)^2\to0$) and the second derivatives ($6(2-x)\to0$, $6(x-2)\to0$) match. So $f_1\in C^2$ and $f_1\in S_{3,\Delta}$; only the third derivative jumps, which is allowed. \pts{1}

$f_2$ is $0$ on $[0,\frac12)$ and $(x-\frac12)^3$ on $[\frac12,1]$, so it is not a single polynomial on $[0,1]$. The point $\frac12$ is not a knot, so $f_2\notin S_{3,\Delta}$, although $f_2\in C^2$. \pts{1}

$f_3$ is a cubic polynomial on each interval, but $f_3''(x)=6x+2$ for $x>1$ and $6x$ for $x<1$: the second derivative jumps by $2$ at $x=1$. So $f_3\notin C^2$ and $f_3\notin S_{3,\Delta}$. \pts{1}

\sol{b} $\dim S_{3,\Delta}=n+k=3+3=6$ (the knots are $x_0,\dots,x_n$ with $n=3$). \pts{1}
Interpolation gives $n+1=4$ conditions, so $2$ degrees of freedom are left. For example the natural boundary conditions $s''(x_0)=s''(x_n)=0$ fix them. \pts{1}

\sol{c} Natural spline: $M_0=M_3=0$. With $h_i=1$ we have $\lambda_i=\mu_i=\frac12$ and
\[
d_i=\frac6{h_i+h_{i+1}}\Big(\frac{f_{i+1}-f_i}{h_{i+1}}-\frac{f_i-f_{i-1}}{h_i}\Big)=3\,(f_{i+1}-2f_i+f_{i-1}).\ptd{1}
\]
So $d_1=3(0-2+0)=-6$ and $d_2=3(1-0+1)=6$. \pts{1}
The equations $\lambda_iM_{i+1}+2M_i+\mu_iM_{i-1}=d_i$ for $i=1,2$ read
\[
2M_1+\tfrac12M_2=-6,\qquad \tfrac12M_1+2M_2=6.\ptd{1}
\]
Adding both gives $\frac52(M_1+M_2)=0$, so $M_2=-M_1$ and $\frac32M_1=-6$. Hence
$M_0=0$, $M_1=-4$, $M_2=4$, $M_3=0$. \pts{1}

\sol{d} On $[0,1]$, $s''$ is linear from $M_0=0$ to $M_1=-4$: $s''(x)=-4x$. Integrating twice, $s(x)=-\frac23x^3+Cx+D$. From $s(0)=0$ we get $D=0$, and from $s(1)=1$ we get $C=\frac53$:
\[
s(x)=\tfrac53x-\tfrac23x^3,\qquad x\in[0,1].\ptd{1}
\]
$s(\frac12)=\frac56-\frac1{12}=\frac34$. \pts{1}

\begin{trap}
\begin{itemize}
\item Break points must be knots: $f_2$ is $C^2$ but not a spline, while $\abs{x-2}^3$ \emph{is} one.
\item The $M_i$ are second derivatives. To get $s$ you still have to integrate and use the values $s(x_i)$.
\item If you do not remember the system: take $s''$ piecewise linear with unknown $M_i$, integrate, and require $s'$ to be continuous at the inner knots. This gives the same equations.
\end{itemize}
\end{trap}
\begin{heart}
$S_{k,\Delta}$: piecewise in $\mathcal P_k$ and globally $C^{k-1}$; $\dim S_{k,\Delta}=n+k$. Natural spline: $s''(x_0)=s''(x_n)=0$. Cubic spline system $\lambda_iM_{i+1}+2M_i+\mu_iM_{i-1}=d_i$ with $\lambda_i=\frac{h_{i+1}}{h_i+h_{i+1}}$, $\mu_i=\frac{h_i}{h_i+h_{i+1}}$.
\end{heart}
\end{solution}

\end{document}
"
\documentclass[11pt,a4paper]{article}
\usepackage[T1]{fontenc}
\usepackage[utf8]{inputenc}
\usepackage{lmodern}
\usepackage{amsmath,amssymb,mathtools,bm}
\usepackage[margin=2.3cm,top=2.2cm,bottom=2.2cm,headsep=0.5cm]{geometry}
\usepackage{enumitem}
\usepackage{xcolor}
\usepackage[most]{tcolorbox}
\usepackage{fancyhdr}
\usepackage{lastpage}
\usepackage{needspace}
\usepackage{comment}
\usepackage{array}

\setlength{\parindent}{0pt}
\setlength{\parskip}{0.4em}
\AtBeginDocument{%
  \setlength{\abovedisplayskip}{6pt plus 2pt minus 2pt}%
  \setlength{\belowdisplayskip}{6pt plus 2pt minus 2pt}%
  \setlength{\abovedisplayshortskip}{3pt plus 2pt}%
  \setlength{\belowdisplayshortskip}{4pt plus 2pt minus 1pt}}

% ---------- notation (follows the formula sheet) ----------
\newcommand{\R}{\mathbb{R}}
\newcommand{\C}{\mathbb{C}}
\newcommand{\Z}{\mathbb{Z}}
\newcommand{\N}{\mathbb{N}}
\newcommand{\T}{^{\mathsf{T}}}
\newcommand{\vb}[1]{\mathbf{#1}}
\newcommand{\ten}[1]{\mathcal{#1}}
\newcommand{\dd}{\,\mathrm{d}}
\newcommand{\e}{\mathrm{e}}
\newcommand{\iu}{\mathrm{i}}
\newcommand{\F}{\mathcal{F}}
\newcommand{\norm}[1]{\lVert #1\rVert}
\newcommand{\abs}[1]{\lvert #1\rvert}
\newcommand{\ip}[2]{\langle #1,#2\rangle}
\newcommand{\op}{\circledcirc}          % outer product as in the lecture
\newcommand{\kr}{\odot}                 % Khatri--Rao
\newcommand{\had}{\ast}                 % Hadamard
\newcommand{\Hess}{\operatorname{Hess}}
\newcommand{\rank}{\operatorname{rank}}
\newcommand{\spn}{\operatorname{span}}
\DeclareMathOperator{\sinc}{sinc}

% ---------- points and boxes ----------
% points at the end of a paragraph, always flush right (TeXbook trick)
\newcommand{\pts}[1]{{\unskip\nobreak\hfil\penalty50\hskip1.5em\hbox{}\nobreak\hfil\textit{(#1\,P.)}\parfillskip=0pt\finalhyphendemerits=0\par}}
% points inside a display, placed like an equation tag
\newcommand{\ptd}[1]{\tag*{\textit{(#1\,P.)}}}
\newcommand{\ptsin}[1]{\textit{(#1\,P.)}}

\newtcolorbox{infobox}{colback=gray!8,colframe=gray!45,arc=2pt,boxrule=0.5pt,left=6pt,right=6pt,top=4pt,bottom=4pt}
\newtcolorbox{trap}{enhanced,colback=red!4,colframe=red!55!black,arc=2pt,boxrule=0.5pt,
  title={\textbf{Typical trap}},fonttitle=\small\bfseries,fontupper=\small,left=6pt,right=6pt,top=2pt,bottom=2pt,before skip=8pt,after skip=6pt}
\newtcolorbox{heart}{enhanced,colback=blue!4,colframe=blue!45!black,arc=2pt,boxrule=0.5pt,
  title={\textbf{Not on the formula sheet -- know by heart}},fonttitle=\small\bfseries,fontupper=\small,left=6pt,right=6pt,top=2pt,bottom=2pt,before skip=8pt,after skip=6pt}

% ---------- problems ----------
\newcounter{problem}
\newcommand{\problem}[2]{%
  \ifnum\value{problem}>0
    \ifdefined\withsol\par\vspace{1.4em}\Needspace{0.18\textheight}{\color{gray!60}\hrule height 0.6pt}\vspace{0.9em}%
    \else\par\bigskip\Needspace{10\baselineskip}\fi
  \fi
  \stepcounter{problem}%
  {\large\textbf{Problem \theproblem.\ (#1, #2 points)}}\par\smallskip}

\setlist[enumerate,1]{label=\textbf{\alph*)},leftmargin=2.2em,itemsep=0.35em,topsep=0.3em}
\setlist[itemize]{itemsep=0.15em,topsep=0.2em}
\newcommand{\sol}[1]{\par\medskip\textbf{#1)}~}

\ifdefined\withsol
  \specialcomment{solution}{\par\medskip\noindent{\large\textbf{Solution.}}\par\nopagebreak}{\par}
\else
  \excludecomment{solution}
\fi

% ---------- header and title ----------
\newcommand{\examletter}{X}
\newcommand{\makeheader}{%
  \pagestyle{fancy}\fancyhf{}
  \renewcommand{\headrulewidth}{0.4pt}
  \fancyhead[L]{\small TMA4432 Mathematics 4C -- Practice exam \examletter\ifdefined\withsol\ (with solutions)\fi}
  \fancyhead[R]{\small Page \thepage\ of \pageref{LastPage}}
  \thispagestyle{fancy}
  \begin{center}
    {\LARGE\textbf{Practice exam \examletter}}\\[4pt]
    {\large TMA4432 Mathematics 4C, autumn 2026\ifdefined\withsol\ -- solutions\fi}
  \end{center}
  \begin{infobox}
  \textbf{Time:} 4 hours. \textbf{Points:} 100. \textbf{Aids (category D):} the official formula sheet and a simple calculator, nothing else.\\
  All answers have to be justified, with enough detail to see how they were obtained.\\[3pt]
  \begin{tabular}{@{}l@{\hspace{1.2em}}*{5}{p{1.4cm}}l@{}}
  Grade & A & B & C & D & E & F\\
  Points & 100--89 & 88--77 & 76--65 & 64--53 & 52--41 & 40 and less
  \end{tabular}\\[3pt]
  {\small Unofficial practice exam, built from the syllabus (lectures 1--12 and the summary lecture), the exercise sheets, the test exam and Bergmann's Math 4N exams 2022--2025. Notation follows the formula sheet.}
  \end{infobox}\par\medskip}

\renewcommand{\examletter}{D}
\begin{document}
\makeheader

% =====================================================================
\problem{Rank-deficient least squares}{10}
Consider
\[
A=\begin{pmatrix}1&1\\2&2\\2&2\end{pmatrix}\in\R^{3\times2},\qquad \vb b=\begin{pmatrix}3\\0\\3\end{pmatrix}.
\]
\begin{enumerate}
\item Determine the rank of $A$ and a reduced singular value decomposition of $A$ (only nonzero singular values).
\item Compute the Moore--Penrose inverse $A^\dagger$.
\item Set up the normal equations for $\min_{\vb x}\norm{A\vb x-\vb b}_2$. Why do they not have a unique solution? Describe all least-squares solutions and state their error $\norm{A\vb x-\vb b}_2$.
\item Compute $\vb x^\dagger=A^\dagger\vb b$ and show that it is the least-squares solution of smallest norm.
\end{enumerate}

\begin{solution}
\sol{a} All rows are multiples of $(1,1)$, so $\rank A=1$. \pts{1}
$A\T A=\begin{psmallmatrix}9&9\\9&9\end{psmallmatrix}$ has the eigenvalues $18$ and $0$, so $\sigma_1=\sqrt{18}=3\sqrt2$ and $\sigma_2=0$. With $\vb v_1=\frac1{\sqrt2}(1,1)\T$ we get $\vb u_1=A\vb v_1/\sigma_1=\frac{1}{3\sqrt2}\cdot\frac{1}{\sqrt2}(2,4,4)\T=\frac13(1,2,2)\T$. \pts{1}
The reduced SVD consists of one term:
\[
A=\sigma_1\vb u_1\vb v_1\T=3\sqrt2\cdot\frac13\begin{pmatrix}1\\2\\2\end{pmatrix}\cdot\frac1{\sqrt2}\begin{pmatrix}1&1\end{pmatrix}.\ptd{1}
\]

\sol{b} Only the nonzero singular value is inverted:
\[
A^\dagger=\frac1{\sigma_1}\vb v_1\vb u_1\T=\frac1{3\sqrt2}\cdot\frac1{\sqrt2}\begin{pmatrix}1\\1\end{pmatrix}\cdot\frac13\begin{pmatrix}1&2&2\end{pmatrix}=\frac1{18}\begin{pmatrix}1&2&2\\1&2&2\end{pmatrix}.\ptd{2}
\]

\sol{c} $A\T A\vb x=A\T\vb b$ reads $\begin{psmallmatrix}9&9\\9&9\end{psmallmatrix}\vb x=\begin{psmallmatrix}9\\9\end{psmallmatrix}$, i.e.\ $x_1+x_2=1$. \pts{1}
$A\T A$ is singular (rank 1), so there are infinitely many least-squares solutions $\vb x=(1-\tau,\tau)\T$, $\tau\in\R$. For all of them $A\vb x=(x_1+x_2)(1,2,2)\T=(1,2,2)\T$, so $A\vb x-\vb b=(-2,2,-1)\T$ and $\norm{A\vb x-\vb b}_2=3$. \pts{1}

\sol{d} $\vb x^\dagger=A^\dagger\vb b=\frac1{18}(3+0+6)(1,1)\T=(\frac12,\frac12)\T$. It satisfies $x_1+x_2=1$, so it is a least-squares solution. \pts{1}
$\ker A=\spn\{(1,-1)\T\}$ and every least-squares solution can be written as $\vb x=\vb x^\dagger+\rho\,(1,-1)\T$ ($\rho=\frac12-\tau$). Since $\ip{\vb x^\dagger}{(1,-1)\T}=0$,
\[
\norm{\vb x}_2^2=\norm{\vb x^\dagger}_2^2+2\rho^2\ \ge\ \tfrac12,
\]
with equality only for $\rho=0$. So $\vb x^\dagger$ is the smallest-norm vector among all vectors with the smallest error: ``smallest norm among those closest''. \pts{2}

\begin{trap}
\begin{itemize}
\item For rank-deficient $A$ the formula $(A\T A)^{-1}A\T$ does not exist. Use the SVD and invert only $\sigma_j>0$.
\item The least-squares solution is not unique here, but the residual is the same for all of them.
\end{itemize}
\end{trap}
\begin{heart}
$A^\dagger=\sum_{\sigma_j>0}\frac1{\sigma_j}\vb v_j\vb u_j\T$. $\vb x^\dagger=A^\dagger\vb b$ minimizes $\norm{\vb x}_2$ among all minimizers of $\norm{A\vb x-\vb b}_2$. Normal equations: $A\T A\vb x=A\T\vb b$.
\end{heart}
\end{solution}

% =====================================================================
\problem{Rank-1 tensors and the Khatri--Rao product}{8}
Let $\vb a=(1,1)\T$, $\vb b=(1,-1)\T$ and
\[
\ten X=\vb a\op\vb a\op\vb a+\vb b\op\vb b\op\vb b\in\R^{2\times2\times2}.
\]
\begin{enumerate}
\item Compute the frontal slices $X_{::1}$ and $X_{::2}$.
\item Write $\ten X=[\![A,B,C]\!]$ with $A=B=C=(\vb a\ \vb b)$. Compute $C\kr B$ and verify $X_{(1)}=A(C\kr B)\T$.
\item Determine the rank of $\ten X$.
\item Compute the 3-mode products $\ten X\times_3\vb w$ for $\vb w=(1,1)\T$ and $\vb w=(1,-1)\T$, and explain the result.
\end{enumerate}

\begin{solution}
\sol{a} $x_{ijk}=a_ia_ja_k+b_ib_jb_k=1+b_ib_jb_k$. For $k=1$ ($b_1=1$): $x_{ij1}=1+b_ib_j$; for $k=2$ ($b_2=-1$): $x_{ij2}=1-b_ib_j$. Hence
\[
X_{::1}=\begin{pmatrix}2&0\\0&2\end{pmatrix},\qquad X_{::2}=\begin{pmatrix}0&2\\2&0\end{pmatrix}.\ptd{2}
\]

\sol{b} The Khatri--Rao product is the columnwise Kronecker product:
\[
C\kr B=\big(\vb a\otimes\vb a\ \ \vb b\otimes\vb b\big)=\begin{pmatrix}1&1\\1&-1\\1&-1\\1&1\end{pmatrix}.\ptd{1}
\]
\[
A(C\kr B)\T=\begin{pmatrix}1&1\\1&-1\end{pmatrix}\begin{pmatrix}1&1&1&1\\1&-1&-1&1\end{pmatrix}=\begin{pmatrix}2&0&0&2\\0&2&2&0\end{pmatrix}=\big(X_{::1}\ X_{::2}\big)=X_{(1)}.\ptd{1}
\]

\sol{c} By construction $\rank\ten X\le2$. A sum of $r$ rank-1 tensors has $\rank X_{(1)}\le r$, and the two rows of $X_{(1)}$ are linearly independent. So $\rank X_{(1)}=2$ and $\rank\ten X=2$. \pts{2}

\sol{d} $\ten X\times_3\vb w=\sum_kx_{ijk}w_k=w_1X_{::1}+w_2X_{::2}$:
\[
\ten X\times_3\begin{pmatrix}1\\1\end{pmatrix}=\begin{pmatrix}2&2\\2&2\end{pmatrix}=2\,\vb a\vb a\T,\qquad
\ten X\times_3\begin{pmatrix}1\\-1\end{pmatrix}=\begin{pmatrix}2&-2\\-2&2\end{pmatrix}=2\,\vb b\vb b\T.\ptd{1}
\]
For a rank-1 tensor, $(\vb a\op\vb a\op\vb a)\times_3\vb w=\ip{\vb a}{\vb w}\vb a\vb a\T$. Hence $\ten X\times_3\vb w=\ip{\vb a}{\vb w}\vb a\vb a\T+\ip{\vb b}{\vb w}\vb b\vb b\T$. Since $\ip{\vb a}{(1,1)\T}=2$ and $\ip{\vb b}{(1,1)\T}=0$, only the first term survives, and vice versa for $(1,-1)\T$. The product with a vector reduces the order by one: the result is a matrix. \pts{1}

\begin{trap}
\begin{itemize}
\item Every entry of $\vb a\op\vb b\op\vb c$ is a product of \emph{three} numbers, $a_ib_jc_k$.
\item In $X_{(1)}=A(C\kr B)\T$ the order is $C\kr B$, not $B\kr C$.
\end{itemize}
\end{trap}
\begin{heart}
$A\kr B=(\vb a_1\otimes\vb b_1\ \cdots\ \vb a_n\otimes\vb b_n)$; $X_{(1)}=A(C\kr B)\T$. Mode product with a vector: $(\ten X\times_k\vb a)=\sum_{i_k}x_{\dots i_k\dots}a_{i_k}$. $\rank\ten X\ge\rank X_{(k)}$.
\end{heart}
\end{solution}

% =====================================================================
\problem{A parametrized surface}{8}
Consider the surface $S$ parametrized by
\[
\vb s(u,v)=\begin{pmatrix}u\cos v\\u\sin v\\u\end{pmatrix},\qquad (u,v)\in D=[0,\infty)\times[0,2\pi).
\]
\begin{enumerate}
\item Show that $S$ lies on the cone $x^2+y^2=z^2$, $z\ge0$. Compute the velocity vectors $\vb s_u,\vb s_v$ and the normal vector $\vb s_u\times\vb s_v$.
\item In which points is the parametrization not smooth?
\item Show that $\vb p=(3,4,5)\T\in S$ and determine the tangent plane in $\vb p$, both in parametric form and as an equation.
\item Compare with the gradient of $F(x,y,z)=x^2+y^2-z^2$. Is there a tangent plane in the apex $(0,0,0)\T$?
\end{enumerate}

\begin{solution}
\sol{a} $(u\cos v)^2+(u\sin v)^2=u^2=z^2$ with $z=u\ge0$. \pts{1}
\[
\vb s_u=\begin{pmatrix}\cos v\\\sin v\\1\end{pmatrix},\qquad \vb s_v=\begin{pmatrix}-u\sin v\\u\cos v\\0\end{pmatrix},
\]
\[
\vb s_u\times\vb s_v=\begin{pmatrix}\sin v\cdot0-1\cdot u\cos v\\1\cdot(-u\sin v)-\cos v\cdot0\\\cos v\cdot u\cos v+\sin v\cdot u\sin v\end{pmatrix}=\begin{pmatrix}-u\cos v\\-u\sin v\\u\end{pmatrix}.\ptd{1}
\]

\sol{b} The total derivative $\vb s'(u,v)=(\vb s_u\ \vb s_v)\in\R^{3\times2}$ has rank 2 iff $\vb s_u\times\vb s_v\neq\vb 0$. Here $\norm{\vb s_u\times\vb s_v}_2=\sqrt2\,u$, which is zero only for $u=0$, where $\vb s_v=\vb 0$. \pts{1}
So $\vb s$ is not smooth at $u=0$, the apex. There it is not injective either: every $v$ is mapped to the origin. \pts{1}

\sol{c} $u=z=5$, $\cos v_0=\frac35$, $\sin v_0=\frac45$ give $\vb s(5,v_0)=(3,4,5)\T$, so $\vb p\in S$. \pts{1}
There $\vb s_u=(\frac35,\frac45,1)\T$ and $\vb s_v=(-4,3,0)\T$, and the tangent plane is $\{\vb p+\sigma\vb s_u+\tau\vb s_v:\sigma,\tau\in\R\}$. \pts{1}
The normal is $\vb n=\vb s_u\times\vb s_v=(-3,-4,5)\T$. The equation is $-3(x-3)-4(y-4)+5(z-5)=0$, i.e.\ $3x+4y-5z=0$. \pts{1}

\sol{d} $\nabla F=(2x,2y,-2z)\T$ and $\nabla F(\vb p)=(6,8,-10)\T=-2\vb n$: the gradient is orthogonal to the level set $F=0$ and gives the same plane. At the apex $\nabla F(\vb 0)=\vb 0$, which gives no normal. The cone has a tip there and no tangent plane: it is only piecewise smooth. \pts{1}

\begin{trap}
\begin{itemize}
\item To find the tangent plane in $\vb p$ you first need the parameters $(u,v)$ with $\vb s(u,v)=\vb p$.
\item The normal is a cross product of the two velocity vectors; the plane passes through $\vb p$.
\end{itemize}
\end{trap}
\begin{heart}
Tangent plane $\vb p+\sigma\vb s_u+\tau\vb s_v$, normal $\vb s_u\times\vb s_v$. Smooth parametrization: injective and $\vb s_u\times\vb s_v\neq\vb 0$ (total derivative of rank 2). $\nabla F$ is orthogonal to the level set $F=c$.
\end{heart}
\end{solution}

% =====================================================================
\problem{Jacobi matrix and Newton's method}{7}
\begin{enumerate}
\item Let $\vb G\colon\R^3\to\R^2$, $\vb G(x,y,z)=\big(xy^2-z,\ \sin(xz)+y\big)\T$. Compute the Jacobi matrix $\vb G'(x,y,z)$ and the linearization $\vb L$ of $\vb G$ in $\vb a=(0,1,2)\T$.
\item Let $\vb F(x,y)=\big(x^2-y^2-3,\ xy-2\big)\T$. Perform one step of Newton's method starting from $(x_0,y_0)=(1,1)$.
\item Where is $\vb F'$ singular? Which zero of $\vb F$ do you expect the iteration from b) to approach?
\end{enumerate}

\begin{solution}
\sol{a} One row per component, one column per variable:
\[
\vb G'(x,y,z)=\begin{pmatrix}y^2&2xy&-1\\z\cos(xz)&1&x\cos(xz)\end{pmatrix}\in\R^{2\times3}.\ptd{1}
\]
$\vb G(\vb a)=(-2,1)\T$ and $\vb G'(\vb a)=\begin{psmallmatrix}1&0&-1\\2&1&0\end{psmallmatrix}$, so
\[
\vb L(x,y,z)=\vb G(\vb a)+\vb G'(\vb a)\begin{pmatrix}x\\y-1\\z-2\end{pmatrix}=\begin{pmatrix}-2+x-(z-2)\\1+2x+(y-1)\end{pmatrix}=\begin{pmatrix}x-z\\2x+y\end{pmatrix}.\ptd{2}
\]

\sol{b} $\vb F(1,1)=(-3,-1)\T$, $\vb F'(x,y)=\begin{psmallmatrix}2x&-2y\\y&x\end{psmallmatrix}$ and $\vb F'(1,1)=\begin{psmallmatrix}2&-2\\1&1\end{psmallmatrix}$ with determinant $4$. \pts{1}
\[
\Delta=-\vb F'(1,1)^{-1}\vb F(1,1)=-\frac14\begin{pmatrix}1&2\\-1&2\end{pmatrix}\begin{pmatrix}-3\\-1\end{pmatrix}=\begin{pmatrix}5/4\\-1/4\end{pmatrix},\qquad (x_1,y_1)=\big(\tfrac94,\tfrac34\big).\ptd{2}
\]

\sol{c} $\det\vb F'(x,y)=2x^2+2y^2=0$ only at $(0,0)$, so Newton can be started anywhere else. The zeros satisfy $x^2-y^2=3$, $xy=2$, i.e.\ $(2,1)$ and $(-2,-1)$. The iteration approaches $(2,1)$: $\vb x^{(2)}\approx(1.992,\,0.975)$, $\vb x^{(3)}\approx(1.99993,\,1.00014)$. \pts{1}

\begin{trap}
\begin{itemize}
\item The Jacobi matrix of $\vb G\colon\R^3\to\R^2$ is $2\times3$: rows are components, columns are variables.
\item The linearization uses $\vb x-\vb a$, not $\vb x$.
\end{itemize}
\end{trap}
\end{solution}

% =====================================================================
\problem{Constrained optimization}{10}
Consider $f(x,y)=xy$ on the ellipse $E=\{(x,y)\T:g(x,y)=0\}$ with $g(x,y)=x^2+4y^2-8$.
\begin{enumerate}
\item Compute $\nabla f$ and $\nabla g$. Compute the directional derivative of $f$ in $\vb p=(2,1)\T$ in the direction of the tangent to $E$ in $\vb p$, and interpret the result.
\item Use Lagrange multipliers to find all candidates for extrema of $f$ on $E$. Determine the maximum and the minimum.
\item Show that the level curve $\{f=2\}$ and the ellipse $E$ have a common tangent line in $\vb p$.
\item State the Lagrangian $\mathcal L(x,y,\lambda)$. Why can we not simply minimize $\mathcal L$ to find the constrained minimum?
\end{enumerate}

\begin{solution}
\sol{a} $\nabla f=(y,x)\T$ and $\nabla g=(2x,8y)\T$. \pts{1}
At $\vb p$ we have $\nabla g(\vb p)=(4,8)\T$, so the tangent direction of $E$ is $\vb t=\frac1{\sqrt5}(2,-1)\T\perp\nabla g(\vb p)$. With $\nabla f(\vb p)=(1,2)\T$:
$D_{\vb t}f(\vb p)=\ip{\nabla f(\vb p)}{\vb t}=\frac{2-2}{\sqrt5}=0$. Moving along the ellipse does not change $f$ to first order in $\vb p$, which is necessary for an extremum under the constraint. \pts{1}

\sol{b} The conditions $\nabla f=\lambda\nabla g$, $g=0$ read
\[
y=2\lambda x,\qquad x=8\lambda y,\qquad x^2+4y^2=8.\ptd{1}
\]
Inserting gives $y=2\lambda\cdot8\lambda y=16\lambda^2y$. If $y=0$, then $x=8\lambda y=0$, but $(0,0)\notin E$. So $16\lambda^2=1$, i.e.\ $\lambda=\pm\frac14$. \pts{1}
$\lambda=\frac14$: $y=\frac x2$, so $x^2+x^2=8$ and $x=\pm2$: the points $(2,1)$, $(-2,-1)$ with $f=2$.
$\lambda=-\frac14$: $y=-\frac x2$: the points $(2,-1)$, $(-2,1)$ with $f=-2$. \pts{2}
$E$ is closed and bounded and $f$ is continuous, so the maximum and the minimum exist; since $\nabla g\neq\vb 0$ on $E$, they are among the candidates. Maximum $f=2$ in $\pm(2,1)$, minimum $f=-2$ in $\pm(2,-1)$. \pts{1}

\sol{c} The level curve $xy=2$ is $y=2/x$ with slope $-2/x^2=-\frac12$ at $x=2$, i.e.\ direction $(2,-1)\T$. This is the tangent direction of $E$ from a). Equivalently, $\nabla f(\vb p)=(1,2)\T=\frac14(4,8)\T=\frac14\nabla g(\vb p)$: the normals are parallel, so the curves touch in $\vb p$. \pts{2}

\sol{d} $\mathcal L(x,y,\lambda)=f(x,y)+\lambda g(x,y)$. $\nabla\mathcal L=\vb 0$ gives $\nabla f+\lambda\nabla g=\vb 0$ and $g=0$: the same points as in b), with $\lambda$ of opposite sign. These points are in general saddle points of $\mathcal L$ (it is linear in $\lambda$), so a minimization method like gradient descent applied to $\mathcal L$ does not find them. \pts{1}

\begin{trap}
\begin{itemize}
\item Do not divide by $y$ (or $x$) without treating the case $y=0$ separately.
\item With $\mathcal L=f+\lambda g$ one gets $\nabla f=-\lambda\nabla g$, the theorem is stated as $\nabla f=\lambda\nabla g$. The sign of $\lambda$ changes, the points do not.
\end{itemize}
\end{trap}
\begin{heart}
If $\vb x^*$ is an extremum of $f$ under $g=0$ and $\nabla g(\vb x^*)\neq\vb 0$, then $\nabla f(\vb x^*)=\lambda\nabla g(\vb x^*)$. Lagrangian $\mathcal L=f+\lambda g$. Directional derivative $D_{\vb v}f=\ip{\nabla f}{\vb v}$, $\norm{\vb v}_2=1$.
\end{heart}
\end{solution}

% =====================================================================
\problem{Fourier transform}{12}
Let $f(x)=1$ for $\abs x<1$ and $f(x)=0$ otherwise.
\begin{enumerate}
\item State the Fourier transform $\hat f(\omega)$. What is $\hat f(0)$?
\item Compute the convolution $g=f\ast f$, $(f\ast f)(x)=\int_{-\infty}^\infty f(y)f(x-y)\dd y$.
\item Use the convolution theorem to compute $\hat g$. Check the value $\hat g(0)$ directly from $g$.
\item Use the rule $\F(h')=\iu\omega\F(h)$ and the formula sheet to compute the Fourier transform of $k(x)=x\,\e^{-x^2}$.
\end{enumerate}

\begin{solution}
\sol{a} The formula sheet with $a=1$ gives $\hat f(\omega)=\sqrt{\frac2\pi}\,\frac{\sin\omega}{\omega}$ for $\omega\neq0$. \pts{1}
$\hat f(0)=\frac1{\sqrt{2\pi}}\int_{-1}^1 1\dd x=\frac2{\sqrt{2\pi}}=\sqrt{\frac2\pi}$, which is also the limit for $\omega\to0$ since $\frac{\sin\omega}\omega\to1$. \pts{1}

\sol{b} $f(x-y)=1$ iff $\abs{x-y}<1$ iff $y\in(x-1,x+1)$. So $g(x)$ is the length of $(-1,1)\cap(x-1,x+1)$. \pts{1}
For $0\le x\le2$ the intersection is $(x-1,1)$, of length $2-x$. For $-2\le x\le0$ it is $(-1,x+1)$, of length $2+x$. For $\abs x\ge2$ it is empty. \pts{2}
Hence $g(x)=2-\abs x$ for $\abs x\le2$ and $g(x)=0$ otherwise, a triangle (hat function). \pts{1}

\sol{c} With $\F(f\ast g)=\sqrt{2\pi}\,\F(f)\F(g)$:
\[
\hat g(\omega)=\sqrt{2\pi}\,\hat f(\omega)^2=\sqrt{2\pi}\cdot\frac2\pi\cdot\frac{\sin^2\omega}{\omega^2}=2\sqrt{\frac2\pi}\,\frac{\sin^2\omega}{\omega^2}.\ptd{2}
\]
Check: $\hat g(0)=\frac1{\sqrt{2\pi}}\int g\dd x=\frac{4}{\sqrt{2\pi}}=2\sqrt{\frac2\pi}$ (the triangle has base $4$ and height $2$, area $4$), the same as the limit of the formula. \pts{1}

\sol{d} Let $h(x)=\e^{-x^2}$. The table gives $\hat h(\omega)=\frac1{\sqrt2}\e^{-\omega^2/4}$, and $h'(x)=-2x\e^{-x^2}=-2k(x)$. \pts{1}
So $-2\hat k(\omega)=\iu\omega\hat h(\omega)$ and
\[
\hat k(\omega)=-\frac{\iu\omega}{2\sqrt2}\,\e^{-\omega^2/4}.\ptd{2}
\]
$k$ is real and odd, so $\hat k$ is purely imaginary and odd, as expected.

\begin{trap}
\begin{itemize}
\item With the normalization of the formula sheet the convolution theorem has the factor $\sqrt{2\pi}$.
\item $\hat f(0)$ has to be taken as a limit (or computed directly), not by inserting $\omega=0$ into $\frac{\sin\omega}{\omega}$.
\end{itemize}
\end{trap}
\begin{heart}
$\F(f\ast g)=\sqrt{2\pi}\,\F(f)\F(g)$, $\F(f')=\iu\omega\F(f)$, $\hat f(0)=\frac1{\sqrt{2\pi}}\int f\dd x$.
\end{heart}
\end{solution}

% =====================================================================
\problem{DFT and aliasing}{10}
Let $f(x)=2\cos x+\e^{3\iu x}$, sampled at $x_j=\frac{2\pi j}4$, $j=0,\dots,3$.
\begin{enumerate}
\item Compute the samples $f_j=f(x_j)$ and the DFT $\hat{\vb f}$ ($N=4$).
\item Explain the result with the aliasing formula. Why is $\hat f_3\neq\overline{\hat f_1}$ here?
\item Which trigonometric polynomial $p(x)=\sum_{k=-1}^{2}\hat f_k\e^{\iu kx}$ (with $\hat f_{-1}=\hat f_3$) is reconstructed from the samples? Compare it with $f$.
\item What is the smallest even $N$ such that the Fourier coefficients of $f$ can be read off from the DFT? Which $N$ would you choose for the FFT?
\end{enumerate}

\begin{solution}
\sol{a} $f_0=2+1=3$, $f_1=2\cos\frac\pi2+\e^{3\pi\iu/2}=-\iu$, $f_2=2\cos\pi+\e^{3\pi\iu}=-3$, $f_3=2\cos\frac{3\pi}2+\e^{9\pi\iu/2}=\iu$. So $\vb f=(3,-\iu,-3,\iu)\T$. \pts{1}
With $w_4=-\iu$ and $\hat f_k=\frac14\sum_jf_j(-\iu)^{jk}$:
\begin{align*}
\hat f_0&=\tfrac14(3-\iu-3+\iu)=0, & \hat f_1&=\tfrac14\big(3+(-\iu)(-\iu)+(-3)(-1)+\iu\cdot\iu\big)=\tfrac14(3-1+3-1)=1,\\
\hat f_2&=\tfrac14(3+\iu-3-\iu)=0, & \hat f_3&=\tfrac14\big(3+(-\iu)\iu+(-3)(-1)+\iu(-\iu)\big)=\tfrac14(3+1+3+1)=2.\ptd{2}
\end{align*}

\sol{b} $f$ has the coefficients $c_1=c_{-1}=1$ (from $2\cos x=\e^{\iu x}+\e^{-\iu x}$) and $c_3=1$. The aliasing formula $\hat f_k=\sum_\ell c_{k+4\ell}$ gives $\hat f_1=c_1=1$ and $\hat f_3=c_3+c_{-1}=2$: the frequency $3$ is aliased onto $-1\equiv3\pmod 4$. Furthermore $\hat f_0=\hat f_2=0$. \pts{2}
$f$ is complex-valued ($\e^{3\iu x}$ is not real), so the symmetry $\hat f_{N-k}=\overline{\hat f_k}$ of real signals does not hold. \pts{1}

\sol{c} $p(x)=\e^{\iu x}+2\e^{-\iu x}$. Since $\e^{3\iu x_j}=\e^{-\iu x_j}$ at all sample points, $p(x_j)=f_j$, but $p\neq f$: with four samples $\e^{3\iu x}$ and $\e^{-\iu x}$ cannot be distinguished. \pts{2}

\sol{d} $f$ contains the frequencies $-1,1,3$. For even $N$ the DFT gives $\hat f_k=c_k$, $k=-\frac N2+1,\dots,\frac N2-1$, if $f$ only contains these frequencies. We need $\frac N2-1\ge3$, i.e.\ $N\ge8$. \pts{1}
$N=8$ is a power of $2$, good for the FFT. Then $\hat f_1=1$, $\hat f_7=\hat f_{-1}=1$, $\hat f_3=1$, exactly $c_1,c_{-1},c_3$. \pts{1}

\begin{trap}
\begin{itemize}
\item Compute the powers $(-\iu)^{jk}$ carefully: $(-\iu)^2=-1$, $(-\iu)^3=\iu$.
\item The conjugate symmetry $\hat f_{N-k}=\overline{\hat f_k}$ only holds for real signals.
\end{itemize}
\end{trap}
\begin{heart}
Aliasing formula $\hat f_k=\sum_{\ell\in\Z}c_{k+\ell N}(f)$. A trigonometric polynomial with highest frequency $m$ needs at least $2m+1$ samples (Nyquist rate).
\end{heart}
\end{solution}

% =====================================================================
\problem{Heat equation with mixed boundary conditions}{14}
Consider
\[
u_t=u_{xx},\quad 0<x<\pi,\ t>0,\qquad u(0,t)=0,\quad u_x(\pi,t)=0,\qquad u(x,0)=3\sin\tfrac x2-\sin\tfrac{5x}2.
\]
\begin{enumerate}
\item Interpret the boundary conditions physically.
\item Use separation of variables $u(x,t)=F(x)G(t)$. Derive the ODEs for $F$ and $G$, and show by treating all three cases of the separation constant that nontrivial solutions satisfying the boundary conditions exist only for $F_n(x)=\sin\big((n-\frac12)x\big)$, $n\in\N$.
\item Determine the solution $u(x,t)$ satisfying the initial condition.
\item What is the temperature at the insulated end $x=\pi$? Which term dominates for large $t$?
\end{enumerate}

\begin{solution}
\sol{a} The left end is kept at temperature $0$ (Dirichlet condition), the right end is insulated: no heat flows through it (Neumann condition). \pts{1}

\sol{b} Inserting gives $F(x)G'(t)=F''(x)G(t)$, i.e.\ $\frac{G'(t)}{G(t)}=\frac{F''(x)}{F(x)}=k$. The left side depends only on $t$, the right side only on $x$, so both are constant. \pts{1}
This gives $F''=kF$ and $G'=kG$. From $u(0,t)=F(0)G(t)=0$ and $u_x(\pi,t)=F'(\pi)G(t)=0$ for all $t$, with $G\not\equiv0$, we get $F(0)=0$ and $F'(\pi)=0$. \pts{1}

$k=0$: $F=A+Bx$, $F(0)=A=0$, $F'(\pi)=B=0$, so $F\equiv0$. \pts{1}

$k=\mu^2>0$: $F=A\e^{\mu x}+B\e^{-\mu x}$, $F(0)=A+B=0$ gives $F=2A\sinh(\mu x)$, and $F'(\pi)=2A\mu\cosh(\mu\pi)=0$ gives $A=0$ (as $\cosh>0$). So $F\equiv0$. \pts{1}

$k=-p^2<0$, $p>0$: $F=A\cos(px)+B\sin(px)$, $F(0)=A=0$, $F'(\pi)=Bp\cos(p\pi)=0$. A nontrivial solution needs $\cos(p\pi)=0$, i.e.\ $p=n-\frac12$, $n\in\N$, and $F_n(x)=\sin\big((n-\frac12)x\big)$. \pts{2}
Then $G_n'=-(n-\frac12)^2G_n$, so $G_n(t)=\e^{-(n-\frac12)^2t}$. \pts{1}

\sol{c} By superposition $u(x,t)=\sum_{n\ge1}B_n\e^{-(n-\frac12)^2t}\sin\big((n-\frac12)x\big)$. At $t=0$ we need $\sum_nB_n\sin\big((n-\frac12)x\big)=3\sin\frac x2-\sin\frac{5x}2$: $n=1$ ($p=\frac12$) gives $B_1=3$, $n=3$ ($p=\frac52$) gives $B_3=-1$, all other $B_n=0$. \pts{2}
\[
u(x,t)=3\,\e^{-t/4}\sin\frac x2-\e^{-25t/4}\sin\frac{5x}2.\ptd{2}
\]

\sol{d} $u(\pi,t)=3\e^{-t/4}\sin\frac\pi2-\e^{-25t/4}\sin\frac{5\pi}2=3\e^{-t/4}-\e^{-25t/4}$. \pts{1}
For large $t$ the first term dominates, since it decays slowest ($\e^{-t/4}$). Everywhere $u\to0$: the heat leaves through the left end. \pts{1}

\begin{trap}
\begin{itemize}
\item The boundary conditions for $F$ follow from $u=FG$ only because $G\not\equiv0$.
\item The Neumann condition leads to $\cos(p\pi)=0$, so the eigenvalues are $(n-\frac12)^2$, not $n^2$.
\item Treat $k=0$ and $k>0$ explicitly; they are part of the points.
\end{itemize}
\end{trap}
\begin{heart}
General solutions of $F''=kF$: $A\e^{\mu x}+B\e^{-\mu x}$ ($k=\mu^2>0$), $A+Bx$ ($k=0$), $A\cos px+B\sin px$ ($k=-p^2<0$). Heat equation: $G'=c^2kG$. Superposition of the product solutions.
\end{heart}
\end{solution}

% =====================================================================
\problem{Finite differences for a boundary value problem}{10}
Consider the stationary heat equation with a constant heat source
\[
-u''(x)=2,\quad 0<x<1,\qquad u'(0)=0,\quad u(1)=0.
\]
Use the grid $x_i=ih$, $h=\frac13$, and approximations $U_i\approx u(x_i)$, $i=0,\dots,3$.
\begin{enumerate}
\item Show with Taylor expansions that $\frac{u(x-h)-2u(x)+u(x+h)}{h^2}=u''(x)+\mathcal O(h^2)$.
\item Set up the finite difference scheme. Use a ghost node $x_{-1}=-h$ and a central difference for the boundary condition $u'(0)=0$. Write the system as $A\vb U=\vb b$ for $\vb U=(U_0,U_1,U_2)\T$.
\item Solve the system.
\item Compare with the exact solution $u(x)=1-x^2$ and explain the result.
\end{enumerate}

\begin{solution}
\sol{a} $u(x\pm h)=u\pm hu'+\frac{h^2}2u''\pm\frac{h^3}6u'''+\frac{h^4}{24}u''''(\xi_\pm)$. Adding both expansions, the odd terms cancel: \pts{1}
\[
\frac{u(x-h)-2u(x)+u(x+h)}{h^2}=u''(x)+\frac{h^2}{24}\big(u''''(\xi_+)+u''''(\xi_-)\big)=u''(x)+\mathcal O(h^2).\ptd{1}
\]

\sol{b} At the nodes $i=0,1,2$ (with $U_3=u(1)=0$): $-\frac{U_{i-1}-2U_i+U_{i+1}}{h^2}=2$. \pts{1}
Boundary condition: $\frac{U_1-U_{-1}}{2h}=u'(0)=0$, so $U_{-1}=U_1$, and the equation at $i=0$ becomes $\frac{2U_0-2U_1}{h^2}=2$. \pts{1}
With $\frac1{h^2}=9$:
\[
9\begin{pmatrix}2&-2&0\\-1&2&-1\\0&-1&2\end{pmatrix}\begin{pmatrix}U_0\\U_1\\U_2\end{pmatrix}=\begin{pmatrix}2\\2\\2\end{pmatrix}.\ptd{1}
\]

\sol{c} The equations are $U_0-U_1=\frac19$, $-U_0+2U_1-U_2=\frac29$, $-U_1+2U_2=\frac29$. \pts{1}
The last gives $U_1=2U_2-\frac29$, the second $U_0=2U_1-U_2-\frac29=3U_2-\frac69$, and the first $3U_2-\frac69-2U_2+\frac29=\frac19$. Hence
\[
U_2=\tfrac59,\qquad U_1=\tfrac89,\qquad U_0=1.\ptd{2}
\]

\sol{d} The exact values are $u(0)=1$, $u(\frac13)=\frac89$, $u(\frac23)=\frac59$: the numerical solution is exact. \pts{1}
The reason: $u$ is a quadratic polynomial, so $u'''=u''''=0$. The remainder of the second difference vanishes, and so does the error $\frac{h^2}6u'''$ of the central difference for $u'(0)$. \pts{1}

\begin{trap}
\begin{itemize}
\item In general the ghost node gives $U_{-1}=U_1-2h\,u'(0)$. Keep the equation in $x_0$: $U_0$ is unknown for a Neumann condition.
\item The Dirichlet value $U_3$ is known and moves to the right-hand side (here it is $0$).
\end{itemize}
\end{trap}
\begin{heart}
Central difference $u'(x)\approx\frac{u(x+h)-u(x-h)}{2h}$, error $\mathcal O(h^2)$; second difference $\frac{u(x-h)-2u(x)+u(x+h)}{h^2}$, error $\mathcal O(h^2)$; ghost nodes for Neumann and Robin conditions.
\end{heart}
\end{solution}

% =====================================================================
\problem{B-splines}{11}
Consider the uniform knots $t_i=i$, $i\in\Z$, the B-splines of order 1, $B_{i,1}(x)=1$ for $x\in[t_i,t_{i+1})$ and $0$ otherwise, and the Cox--de Boor recursion
\[
B_{i,k+1}(x)=\frac{x-t_i}{t_{i+k}-t_i}B_{i,k}(x)+\frac{t_{i+k+1}-x}{t_{i+k+1}-t_{i+1}}B_{i+1,k}(x).
\]
\begin{enumerate}
\item Compute $B_{0,2}$, $B_{1,2}$ and the quadratic B-spline $B_{0,3}$ piecewise.
\item Show that $B_{0,3}\in C^1(\R)$ and verify the partition of unity $B_{-2,3}+B_{-1,3}+B_{0,3}=1$ on $[0,1)$. You may use $B_{i,3}(x)=B_{0,3}(x-i)$.
\item Let $\Delta=[0,1,2]$. Determine $\alpha,\beta,\gamma$ such that
\[
f(x)=\begin{cases}x^3,&x\in[0,1),\\-x^3+\alpha x^2+\beta x+\gamma,&x\in[1,2],\end{cases}
\]
is in $S_{3,\Delta}$.
\item State $\dim S_{3,\Delta}$ for $\Delta=[0,1,2]$ and a basis of polynomials and truncated powers.
\end{enumerate}

\begin{solution}
\sol{a} $B_{0,2}(x)=x$ on $[0,1)$, $2-x$ on $[1,2)$; $B_{1,2}(x)=x-1$ on $[1,2)$, $3-x$ on $[2,3)$; both $0$ otherwise. \pts{1}
$B_{0,3}(x)=\frac x2B_{0,2}(x)+\frac{3-x}2B_{1,2}(x)$, so
\[
B_{0,3}(x)=\begin{cases}\frac{x^2}2,&x\in[0,1),\\[2pt]\frac{x(2-x)+(3-x)(x-1)}2=\frac{-2x^2+6x-3}2,&x\in[1,2),\\[2pt]\frac{(3-x)^2}2,&x\in[2,3),\\[2pt]0&\text{otherwise.}\end{cases}\ptd{3}
\]

\sol{b} Values: $\frac12$ from both sides at $x=1$ and $x=2$, $0$ at $x=0$ and $x=3$. Derivatives: at $x=1$, $x\to1$ and $-2x+3\to1$; at $x=2$, $-2x+3\to-1$ and $-(3-x)\to-1$; at $x=0$ and $x=3$ both sides are $0$. So $B_{0,3}\in C^1(\R)$. \pts{1}
On $[0,1)$: $B_{0,3}(x)=\frac{x^2}2$, $B_{-1,3}(x)=B_{0,3}(x+1)=\frac{-2x^2+2x+1}2$ (middle piece) and $B_{-2,3}(x)=B_{0,3}(x+2)=\frac{(1-x)^2}2$ (last piece). The sum is $\frac{x^2-2x^2+2x+1+1-2x+x^2}2=1$. \pts{1}

\sol{c} Both pieces are cubic polynomials. $f\in C^2$ requires equal value, first and second derivative at $x=1$:
\[
-1+\alpha+\beta+\gamma=1,\qquad -3+2\alpha+\beta=3,\qquad -6+2\alpha=6,\ptd{1}
\]
so $\alpha=6$, $\beta=-6$, $\gamma=2$. \pts{2}

\sol{d} $\dim S_{3,\Delta}=n+k=2+3=5$, with the basis $1,\ x,\ x^2,\ x^3,\ (x-1)_+^3$. \pts{2}

\begin{trap}
\begin{itemize}
\item Cox--de Boor uses the \emph{current} order $k$ in $t_{i+k}-t_i$ and $t_{i+k+1}-t_{i+1}$; $B_{i,k+1}$ has degree $k$ and support $[t_i,t_{i+k+1}]$.
\item In c) the cubic coefficient is fixed, so the three conditions determine $\alpha,\beta,\gamma$ uniquely; a general cubic on $[1,2]$ would leave the $x^3$ coefficient free.
\end{itemize}
\end{trap}
\begin{heart}
Cox--de Boor recursion; $\operatorname{supp}B_{i,k+1}=[x_i,x_{i+k+1}]$; partition of unity $\sum_iB_{i,k+1}=1$; $\dim S_{k,\Delta}=n+k$.
\end{heart}
\end{solution}

\end{document}
